

\documentclass{article}
\usepackage{aai2030}
\usepackage{cancel}
\input{macro}

\newcommand{\mn}[1]{\mathnormal{#1}}

\newcommand\ti[1]{\noindent{$\tilde{#1}$}}
\newcommand\hh[1]{\noindent{$\hat{#1}$}}
\newcommand\mm[1]{\noindent{${{#1}}$}}

\begin{document}
\lecturetitle{Optimization}


%%% Linear Algebra Review
%%%%%%%%%%%%%%%%%%%%%%%%%%%%%%%%%%%%%%%%%%%%%%%%%%%%%%%%%%%%%%%%%%%%%%%
\section{Linear Algebra Review} 

\subsection{Rank and Range} \label{subsec:rank}
Let $A\in\fR^{n\times m}$ be a $n\times m$ matrix,
\begin{align}
    A = \begin{bmatrix} a_1&a_2&\dots &a_m\end{bmatrix}
\end{align}
where $a_i\in\fR^n$.
\begin{definition}
The range of $A$ is a column space, the span of all its column vector.
Recall that the span is a set of all possible linear combination of given set of vectors.
\begin{align}
    \mbox{range}(A) = \{x = \sum_{i=1}^m \alpha_i a_i : \alpha_i\in\fR\}\in\fR^n.
\end{align}
\end{definition}
The column space is a vector subspace.
The dimension of vector subspace is defined by the size of basis (the maximum number of linearly independent vector).
\begin{definition}
The rank of matrix $A$ is a dimension of $\mbox{range}(A)$.
\begin{align}
    \mbox{rank}(A) = \mbox{dim}(\mbox{range}(A)).
\end{align}
\end{definition}
Note that the dimension of a row space is also the same as the rank.

\begin{definition}
    The null space of $A$ is a set of vectors that satisfies $Ax = 0$, i.e.,
    \begin{align}
        \mbox{null}(A) = \{x \in\fR^{m}: Ax = 0\}.
    \end{align}
\end{definition}

\begin{theorem}[Rank-Nullity Theorem]
    \begin{align}
        \mbox{rank}(A) + \mbox{dim}(\mbox{null}(A)) = n. 
    \end{align}
\end{theorem}

\begin{definition}
    Orthogonal complement of set $W\subseteq\fR^n$ is given by
    \begin{align}
        W^\perp = \{v\in\fR^n: v^T w = 0\}.
    \end{align}
\end{definition}

\begin{theorem}
    If $W$ is a subspace of $\fR^m$, then
    \begin{align}
        \fR^m = W\oplus W^\perp
    \end{align}
    In other words, for any $x\in \fR^n$, there exists $x_1\in W$ and $x_2\in W^\perp$ such that
    \begin{align}
        x = x_1 + x_2.
    \end{align}
\end{theorem}   

\begin{theorem}
For $A\in\fR^{n\times m}$, the orthogonal complement of row space is a null space of $A$.
The orthogonal complement of a column space of $A$ is a null space of $A^T$.
\end{theorem}



\subsection{Norm and Orthogonal Matrix}\label{subsec:orthogonal}
\begin{definition}
    We say $\|\cdot\|$ is a norm on $\fR^n$ if $\|x\|\in[0, \infty)$ for all $x\in\fR^n$, and
    \begin{itemize}
        \item $\|x+y\| \leq \|x\|+\|y\|$ for all $x, y\in\fR^n$
        \item $\|\alpha x\| = |\alpha|\|x\|$ for all $\alpha\in\fR$ and $x\in\fR^n$
        \item $\|x\|=0$ if and only if $x=0\in\fR^n$
    \end{itemize}
\end{definition}
The Euclidean norm ($\ell_2$-norm) is
\begin{align}
    \|x\|_2 = \sqrt{x^T x} = \sqrt{x_1^2+\dots x_n^2}.
\end{align}
We can generalize $\ell_p$ norm for $p>0$:
\begin{align}
    \|x\|_p = (|x_1|^p+\dots |x_n|^p)^{1/p}.
\end{align}


\begin{definition}
    Let $U\in\fR^{n\times n}$ is orthogonal (or unitary) if any of the following equivalent conditions hold:
    \begin{itemize}
        \item $U^TU=I$
        \item $UU^T=I$
        \item All columns of $U$ forms an orthonormal basis of $\fR^n$
        \item All rows of $U$ forms an orthonormal basis of $\fR^n$
    \end{itemize}
\end{definition}
If $U\in\fR^{n\times n}$ is an orthogonal matrix, then $\|Ux\|_2 = \|x\|_2$ for all $x\in\fR^n$.

\subsection{Singular Value Decomposition}\label{subsec:SVD}
Let $A\in\fR^{m\times n}$. The singular value decomposition (SVD) of $A$ has the form
\begin{align}
A =& U\Sigma V^T\\
=&\begin{bmatrix}u_1& \dots & u_m\end{bmatrix}
 \begin{bmatrix}
 \sigma_1 & 0 & 0 & \dots & 0 & \dots & 0\\
        0 & \sigma_2 & 0 & \dots & 0 & \dots & 0\\
        0 & 0 & \sigma_3 & \dots & 0 & \dots & 0\\
        \vdots &&&&&&\\
        0 & 0 & 0 & \dots & \sigma_r& 0 \dots & 0\\
        0 & 0 & 0 & \dots & 0  & 0 \dots & 0\\
        \vdots &&&&&&\\
        0 & 0 & 0 & \dots & 0  & 0 \dots & 0
 \end{bmatrix}
  \begin{bmatrix}v_1^T \\ v_2^T\\\vdots \\v_n^T\end{bmatrix}\\
  =& \sum_{i=1}^r \sigma_i u_i v_i^T.
\end{align}
where $U\in\fR^{m\times m}, \Sigma\in\fR^{m\times n}, V\in\fR^{n\times n}$, and $r$ is the rank of $A$.
$u_1, \dots, u_m\in\fR^m$ are orthonormal basis of $\fR^m$ and $v_1, \dots, v_n\in\fR^n$ are orthonormal basis of $\fR^n$.
The SVD always exists and can be computed.
Note that we have compact version of SVD:
\begin{align}
A =& U_c \Sigma_c V_c^T\\
=&\underbrace{\begin{bmatrix}u_1& \dots & u_r\end{bmatrix}}_{\in\fR^{m\times r}}
\underbrace{
 \begin{bmatrix}
 \sigma_1 & 0 & 0 & \dots & 0 \\
        0 & \sigma_2 & 0 & \dots & 0\\
        0 & 0 & \sigma_3 & \dots & 0\\
        \vdots &&&&\\
        0 & 0 & 0 & \dots & \sigma_r\\
 \end{bmatrix}}_{\in\fR^{r\times r}}
\underbrace{\begin{bmatrix}v_1^T \\ v_2^T\\\vdots \\v_r^T\end{bmatrix}}_{\in\fR^{r\times n}}
\end{align}
Note that $U_c$ and $V_c$ contain orthonormal columns, but the matrices are not orthogonal matrices:
\begin{align}
    U_c^TU_c = I &\mbox{ but }    U_cU_c^T \neq I \\
    V_c^TV_c = I &\mbox{ but }    V_cV_c^T \neq I .
\end{align}
Also, it is a convention that the singular values are in decreasing order: $\sigma_1\geq \sigma_2\geq\dots \geq\sigma_r\geq0$.

\begin{definition}
    The (Moore--Penrose) pseudo-inverse is as defined as follows:
    \begin{align}
        A^{\dagger} =& V_c \Sigma_c^{-1} U_c^T\\
=&\begin{bmatrix}v_1& \dots & v_r\end{bmatrix}
 \begin{bmatrix}
 \sigma_1^{-1} & 0 & 0 & \dots & 0 \\
        0 & \sigma_2^{-1} & 0 & \dots & 0\\
        0 & 0 & \sigma_3^{-1} & \dots & 0\\
        \vdots &&&&\\
        0 & 0 & 0 & \dots & \sigma_r^{-1}\\
 \end{bmatrix}
\begin{bmatrix}u_1^T \\ u_2^T\\\vdots \\u_r^T\end{bmatrix}
    \end{align}
\end{definition}



%%%%%%%%%%%%%%%%%%%%%%%%%%%%%%%%%%%%%%%%%%%%%%%%%%%%%%%%%%%%%%%%%
%%% Matrix Norms
%%%%%%%%%%%%%%%%%%%%%%%%%%%%%%%%%%%%%%%%%%%%%%%%%%%%%%%%%%%%%%%%%
\section{Matrix Norms, Nuclear Norm, and Low-Rank Approximation}

\subsection{Matrix Norms}

\paragraph{Vector $p$-norm.}
For a vector $x\in\mathbb{R}^n$, the $\ell_p$-norm is
\[
\|x\|_p = \left( \sum_{i=1}^n |x_i|^p \right)^{1/p}, \qquad
\|x\|_\infty = \max_i |x_i|.
\]

\paragraph{Matrix induced $p$-norm.}
For a matrix $A\in\mathbb{R}^{m\times n}$, the induced (or operator) $p$-norm is defined by
\[
\|A\|_p = \max_{x\neq 0} \frac{\|Ax\|_p}{\|x\|_p}.
\]
This definition ensures the inequality
\[
\|Ax\|_p \le \|A\|_p \, \|x\|_p \qquad \text{for all } x,
\]
an essential property of operator norms.

\subsubsection{Spectral Norm ($p=2$)}

For $p=2$, the induced matrix norm is the \emph{spectral norm}:
\[
\|A\|_2 = \max_{x\neq 0} \frac{\|Ax\|_2}{\|x\|_2}.
\]
Let $A=U\Sigma V^\top$ be the singular value decomposition (SVD) of $A$, with singular values
$\sigma_1 \ge \sigma_2 \ge \cdots \ge 0$. Then
\[
\|A\|_2 = \sigma_1.
\]

\paragraph{Why this definition is natural.}
For any matrix norm intended to measure ``how much $A$ can stretch a vector,'' it is essential that
\[
\|Ax\|_2 \le \|A\|_2\|x\|_2
\]
holds for every $x$.  
If we defined the norm in any other arbitrary way, it might fail to satisfy this inequality, making it useless for conditioning analysis, stability bounds, and perturbation theory.  
The quantity
\[
\sup_{x\neq 0} \frac{\|Ax\|_2}{\|x\|_2}
\]
is exactly the minimum number that makes the inequality hold for all $x$, hence giving a principled and mathematically tight definition of the operator norm.

\subsection{Nuclear Norm (Trace Norm)}

Let $A$ have singular values $\sigma_1,\ldots,\sigma_r$.  
The \emph{nuclear norm} (also called the \emph{trace norm}) is defined by
\[
\|A\|_* = \sum_{i=1}^r \sigma_i.
\]

\paragraph{Interpretation as a measure of low rank.}
The nuclear norm is the convex envelope of the matrix rank on the unit spectral-norm ball:
\[
\text{rank}(A) \approx \|A\|_* \quad \text{(convex surrogate)}.
\]
Low-rank matrices have rapidly decaying singular values, making $\|A\|_*$ small.  
Thus the nuclear norm encourages low-rank solutions in convex optimization problems.

\subsection{Matrix Completion and the Netflix Problem}

In the Netflix challenge, the data consists of a partially observed user--movie rating matrix.  
Let
\[
\Omega \subseteq \{1,\ldots,m\}\times\{1,\ldots,n\}
\]
denote the set of observed (user, movie) pairs, and let $M\in\mathbb{R}^{m\times n}$ be the unknown full rating matrix.  
The task is to infer the missing entries assuming that $M$ is well-approximated by a low-rank structure (latent user and movie factors).

\subsubsection{Ideal Low-Rank Formulation}

A natural formulation is to search for a rank-$k$ matrix that best fits the observed entries:
\[
\begin{aligned}
\min_{X} \quad &
\sum_{(i,j)\in\Omega} (X_{ij} - M_{ij})^2 \\
\text{s.t.}\quad &
\operatorname{rank}(X) \le k.
\end{aligned}
\]
This is a direct analogue of the classical low-rank factor model used in collaborative filtering.  
However, the rank constraint makes the problem nonconvex and, in general, NP-hard.

\subsubsection{Convex Relaxation via Nuclear Norm}

To obtain a tractable optimization problem, the rank constraint is replaced with its tightest convex surrogate, the nuclear norm:
\[
\operatorname{rank}(X) 
\;\;\longrightarrow\;\; \|X\|_* = \sum_i \sigma_i(X),
\]
leading to the convex program
\[
\begin{aligned}
\min_{X} \quad &
\sum_{(i,j)\in\Omega} (X_{ij} - M_{ij})^2
\;+\; \lambda \|X\|_* .
\end{aligned}
\]
This relaxation enables efficient algorithms such as singular value thresholding, and under suitable incoherence and sampling assumptions, it can recover the true low-rank matrix exactly.



\subsection{Best Low-Rank Approximation Theorem}

Let $A=U\Sigma V^\top$ be the SVD of $A$, and let
\[
A_k = U \Sigma_k V^\top
\]
denote the rank-$k$ truncation keeping only the top $k$ singular values.

\begin{theorem}[Eckart--Young--Mirsky]
For any rank-$k$ matrix $B$,
\[
\|A - A_k\|_2 = \min_{\operatorname{rank}(B)=k} \|A - B\|_2 = \sigma_{k+1},
\]
\[
\|A - A_k\|_F = \min_{\operatorname{rank}(B)=k} \|A - B\|_F
 = \left( \sum_{i>k} \sigma_i^2 \right)^{1/2}.
\]
Thus $A_k$ is the unique best rank-$k$ approximation of $A$ in both spectral and Frobenius norms.
\end{theorem}

\subsection{Proof of the Best Rank-$k$ Approximation Theorem (2-Norm Case)}

We give a concise proof following the geometric argument in the lecture notes~\footnote{\url{https://rich-d-wilkinson.github.io/MATH3030/3.5-lowrank.html}}..

Let $A \in \mathbb{R}^{n \times p}$ have SVD
\[
A = U \Sigma V^\top,
\qquad
\Sigma = \operatorname{diag}(\sigma_1,\dots,\sigma_p),
\qquad
\sigma_1 \ge \cdots \ge \sigma_p \ge 0.
\]
Let $B$ be any matrix with $\operatorname{rank}(B)=k$.  
We aim to show
\[
\|A - B\|_2 \;\ge\; \sigma_{k+1},
\qquad\text{with equality achieved by } A_k.
\]

\paragraph{Step 1: Use the null space of $B$.}
Since $\operatorname{rank}(B) = k$, the rank–nullity theorem implies
\[
\dim N(B) = p - k.
\]

\paragraph{Step 2: Use the span of the first $k{+}1$ right-singular vectors.}
Consider the matrix
\[
V_{k+1} = [v_1 \; \cdots \; v_{k+1}] \in \mathbb{R}^{p \times (k+1)},
\]
whose columns are the first $k+1$ right singular vectors of $A$.  
This matrix has rank $k{+}1$, so its column space $C(V_{k+1})$ is $(k{+}1)$-dimensional.

\paragraph{Step 3: Dimension argument to find a shared vector.}
Both $N(B)$ and $C(V_{k+1})$ lie in the $p$-dimensional space $\mathbb{R}^p$.  
But
\[
\dim N(B) + \dim C(V_{k+1})
= (p-k) + (k+1)
= p + 1 > p.
\]
Thus they must intersect:
\[
N(B) \cap C(V_{k+1}) \neq \{0\}.
\]
Choose any unit vector $w$ in this intersection:
\[
w \in N(B), \qquad w \in C(V_{k+1}), \qquad \|w\|_2 = 1.
\]

Since $w \in C(V_{k+1})$, we can write
\[
w = \sum_{i=1}^{k+1} w_i v_i,
\qquad
\sum_{i=1}^{k+1} w_i^2 = 1.
\]

\paragraph{Step 4: Lower bound the error using the chosen vector.}
Using the definition of the spectral norm,
\[
\|A - B\|_2^2
\;\ge\;
\|(A - B) w\|_2^2.
\]
But $w \in N(B)$, so $Bw = 0$. Hence
\[
\|(A - B) w\|_2^2
= \|A w\|_2^2
= w^\top V \Sigma^2 V^\top w.
\]
Substituting the expansion of $w$ in the $v_i$ basis,
\[
w^\top V \Sigma^2 V^\top w
= \sum_{i=1}^{k+1} \sigma_i^2 w_i^2
\;\ge\;
\sigma_{k+1}^2 \sum_{i=1}^{k+1} w_i^2
= \sigma_{k+1}^2.
\]
Thus
\[
\|A - B\|_2^2 \;\ge\; \sigma_{k+1}^2,
\qquad
\text{i.e.,} \quad
\|A - B\|_2 \ge \sigma_{k+1}.
\]

\paragraph{Step 5: Tightness.}
For the truncated SVD matrix
\[
A_k = U
\begin{bmatrix}
\operatorname{diag}(\sigma_1,\dots,\sigma_k) & 0 \\
0 & 0
\end{bmatrix}
V^\top,
\]
the residual is
\[
A - A_k = U
\begin{bmatrix}
0 & 0 \\
0 & \operatorname{diag}(\sigma_{k+1},\dots,\sigma_p)
\end{bmatrix}
V^\top,
\]
whose largest singular value is exactly $\sigma_{k+1}$.

Thus for every rank-$k$ matrix $B$,
\[
\|A-A_k\|_2 = \sigma_{k+1} \le \|A-B\|_2.
\]
This completes the proof.
\hfill $\square$


\paragraph{Frobenius norm case.}
Let $B$ be any rank-$k$ matrix.  
Write $A - B$ in the singular basis of $A$:
\[
A = \sum_{i=1}^r \sigma_i u_i v_i^\top, \qquad
A_k = \sum_{i\le k} \sigma_i u_i v_i^\top.
\]
Since $\operatorname{rank}(B)\le k$, it can match at most $k$ singular directions.  
Thus the residual Frobenius energy is at least the energy of the truncated tail:
\[
\|A-B\|_F^2 \ge \sum_{i>k} \sigma_i^2.
\]
Again $A_k$ achieves equality.



%%%%%%%%%%%%%%%%%%%%%%%%%%%%%%%%%%%%%%%%%%%%%%%%%%%%%%%%%%%%%%%%%
%%% Linear Regression
%%%%%%%%%%%%%%%%%%%%%%%%%%%%%%%%%%%%%%%%%%%%%%%%%%%%%%%%%%%%%%%%%
\section{Linear Regression}\label{subsec:linear regression}
Given a data
\begin{align}
    (\bx^{(1)}, y^{(1)}), (\bx^{(2)}, y^{(2)}), \dots,  (\bx^{(N)}, y^{(N)})
\end{align}
where $\bx^{(i)}\in\ \fR^{n}$ and $y^{(i)}\in\fR$,
we want to find a linear model that approximate 
\begin{align}
    f_{\theta}(\bx) = \theta^T \bx = \theta_1x_1 + \dots +\theta_n x_n \approx y
\end{align}
for new data-label pair $(\bx, y)$.

\begin{example}
Suppose $n = 1$, where $x \in \mathbb{R}$ represents an individual's height, and $y$ represents their weight.
To model the relationship between height and weight linearly, we can define the model as:
\begin{align}
    y = \theta_1 x.
\end{align}
Using this model, we can predict an individual's weight from their height by multiplying $\theta_1$ with $x$.
\end{example}

\begin{remark}
The linear model above does not include a bias term.
However, it can be made more general by adding a dummy feature to account for the bias.
Let $\tilde{\mathbf{x}} = (1, x_1, x_2, \dots, x_n)$, and formulate the model to find 
$\tilde{\theta} = (\theta_0, \theta_1, \dots, \theta_n)$ that fits:
\begin{align}
y \approx \tilde{\theta}^T \tilde{\mathbf{x}} = \theta_0 + \theta_1 x_1 + \dots + \theta_n x_n.
\end{align}
This formulation is equivalent to a linear model with a bias term.
\end{remark}

Linear regression is interpretable.
\begin{itemize}
    \item if $\theta_k$ is small, than $x_k$ has less influence on prediction $f_\theta(\bx)$.
    \item if $\theta_k=0$ means the prediction does not depend on $x_k$.
    \item $\theta_k$ is the amount prediction $f_\theta(\bx)$ increases when $x_k$ increases by 1.
\end{itemize}

We need to formally define the goodness (how well approximate $y$).
In linear regression problem, we normally assumes to minimize Mean Squared Error (MSE) or equivalently $\ell_2$-loss.
The goal is to minimizes
\begin{align}
    \ell_2(f_\theta(\bx), y) = (f_\theta(\bx)-y)^2
\end{align}
for new data-label pair $(\bx, y)$.
Given a data, we minimizes the MSE on given dataset.
Then, the problem is to find $\theta^\star$
\begin{align}
    \theta^\star 
    =& \argmin_\theta \frac{1}{N} \sum_{i=1}^N \ell_2(f_\theta(\bx^{(i)}) , y^{(i)})\\
    =& \argmin_\theta \frac{1}{N} \sum_{i=1}^N (f_\theta(\bx^{(i)}) - y^{(i)})^2\\
    =& \argmin_\theta \frac{1}{N} \sum_{i=1}^N (\theta^T\bx^{(i)} - y^{(i)})^2
\end{align}
Note that $\frac{1}{N}$ does not affect the solution.

We can further simplify the above summation by defining
\begin{align}
    A = \begin{bmatrix}
        (\bx^{(1)})^T\\
        (\bx^{(2)})^T\\
        \vdots\\
        (\bx^{(N)})^T\\
    \end{bmatrix}\quad 
    B = \begin{bmatrix}y{(1)}\\y^{(2)}\\\vdots\\y^{(N)}\end{bmatrix}
\end{align}
where $A\in\fR^{N\times n}$ and $B\in\fR^N$.
We normally assume $N > n$ and the matrix $A$ is tall, which means $A\theta = B$ is over-determined (have more equations than variables).
For most cases, there is no exact solution $\theta$ such that $A\theta=B$.

\begin{remark}
    Suppose we are interested in solving linear equation $A\theta = B$. 
    If $A$ is a square invertible matrix, then the solution is $\theta = A^{-1} B$.
    If $A$ is a tall matrix, then we have more equations than the variables, which is over-determined setting,
    which usually means no solution satisfies $A\theta=B$.
    If $A$ is a fat matrix, then we have less equations than the variables, which is under-determined setting,
    which usually have more than 1 solutions that satisfies $Ax=B$ (if $A$ is full rank).
\end{remark}

Then, we have
\begin{align}
    \theta^\star 
    =& \argmin_\theta \frac{1}{2} \|A\theta - B\|_2^2
\end{align}
Again $\frac{1}{2}$ does not affect the solution

\begin{theorem}
    The linear least squares optimization problem
    \begin{align}
        \min_{\theta\in\fR^n} \frac{1}{2}\|A\theta - B\|_2^2.
    \end{align}
    Then, the optimal solution is given by
    \begin{align}
        \theta^\star = A^\dagger B
    \end{align}
    The solution is unique if $n = r\leq N$.
    Otherwise, $\theta^\star$ is not unique but it is the least-norm solution
    that achieves minium $\frac{1}{2}\|A\theta-B\|_2^2$ while having smallest $\|\theta\|_2^2$.
\end{theorem}

\begin{proof}
    Let $A = U_c\Sigma_c V_c^T$ be the compact SVD where $U_c\in\fR^{N\times r}, \Sigma_c\in\fR^{r\times r}, V^T\in \fR^{r\times n}$.
    Let $\tilde{U}\in\fR^{N\times(N-r)}$ be the matrix of $(N-r)$ column vectors so that $U = [U_c~\tilde{U}]\in\fR^{n\times n}$ forms an orthogonal matrix.
    Then,
    \begin{align}
        \argmin_{\theta\in\fR^n} \frac{1}{2}\|A\theta - B\|_2^2 
        &=  \argmin_{\theta\in\fR^n} \frac{1}{2}\|U_c \Sigma_c V_c^T \theta - B\|_2^2\\
        &=  \argmin_{\theta\in\fR^n} \frac{1}{2}\|\Sigma_c V_c^T \theta - U_c^T B\|_2^2 + \frac{1}{2} \|\tilde{U}^TB\|^2\\
        &=  \argmin_{\theta\in\fR^n} \frac{1}{2}\|\Sigma_c V_c^T \theta - U_c^T B\|_2^2
    \end{align}
    In turn, this is equivalent to 
    \begin{align}
        \argmin_{\theta\in\fR^n} \frac{1}{2}\|A\theta - B\|_2^2 
        &=  \argmin_{\theta_1\in\mbox{null}(V_c^T)^\perp, \theta_2\in\mbox{null}(V_c^T)} \frac{1}{2}\|\Sigma_c V_c^T (\theta_1+\theta_2) - U_c^T B\|_2^2\\
        &=  \argmin_{\theta_1\in\mbox{null}(V_c^T)^\perp, \theta_2\in\mbox{null}(V_c^T)} \frac{1}{2}\|\Sigma_c V_c^T \theta_1 - U_c^T B\|_2^2\\
    \end{align}
    Since $V_c^TV_c = I$, the above minimum is achieved by
    \begin{align}
        \theta_1^\star =& V_c\Sigma_c^{-1}U_c^T B\\
        =& A^\dagger B.
    \end{align}
    Thus, $\theta^\star = \theta_1^\star + \theta_2$ is a solution that achieves the minimum for all $\theta_2\in\mbox{null}(V_c^T)$.
    However, among them, the minimum norm solution is $\theta^\star = \theta_1^\star$ since
    \begin{align}
        \|\theta_1^\star + \theta_2\|_2^2 =& \|\theta_1^\star\|_2^2 + \|\theta_2\|_2^2\\
        \geq& \|\theta_1^\star\|_2^2.
    \end{align}
    We will see later why the minimum norm is a favorable property.
        
\end{proof}

\begin{remark}
    It is a common setup in Machine Learning.
    Given a dataset, we first define a function class $\cF$ (in above case, $\cF = \{f_{\theta}:f_{\theta}(\bx) = \theta^T \bx, \theta\in\fR^n\}$),
    and define the loss function $\ell(f_\theta(\bx), y)$.
    Then, minimizes the ``empirical loss'':
    \begin{align}
        \cL (\theta) = \frac{1}{N} \sum_{i=1}^N \ell(f_{\theta}(\bx^{(i)}, y^{(i)}).
    \end{align}
    This is also the case in neural network training where on can set $\cF$ be the class of neural network 
    with appropriate loss $\ell$ (like cross-entropy loss).
    The training is simply to minimize the empirical loss $\cL(\theta)$ with an appropriate optimization algorithm.
\end{remark}

\begin{remark}
We can generalize this approach to a broader class of functions.
For instance, in the $n = 1$ case, we might consider a more complex model such as
$f_\theta(x) = \theta_1 \log x + \theta_2 x + \theta_3 x^2$.
Although this model incorporates nonlinear transformations of $x$, it remains linear with respect to $\theta$, 
so we can apply the same method to find the best model fit.

This approach introduces the problem of selecting the most useful features (like $\log x$, $x^2$, etc.),
a process known as feature engineering. 
In later courses, we will explore feature selection techniques in more depth.
In the deep learning era, however, feature engineering is largely automated,
with neural networks learning the relevant features directly from data.
\end{remark}


More generally, the output dimension can be larger than 1.
\begin{align}
    (\bx^{(1)}, \by^{(1)}), (\bx^{(2)}, \by^{(2)}), \dots,  (\bx^{(N)}, \by^{(N)})
\end{align}
where $\bx^{(i)}\in\ \fR^{n}$ and $\by^{i}\in\fR^m$.

We want to find a linear model that approximate 
\begin{align}
    f_{\Theta}(\bx) = \Theta \bx\approx \by
\end{align}
for new data-label pair $(\bx, \by)$ where $\Theta$ is $n\times m$ matrices.

If we let
\begin{align}
    \Theta = \begin{bmatrix} \theta_1^T \\ \theta_2^T\\ \vdots \\ \theta_m^T\end{bmatrix}.
\end{align}
Then, the problem is simply to find $\theta_k\in\fR^n$ that satisfies $\theta_k \bx \approx \by_k$ for all $1\leq k \leq m$.
This is because
\begin{align}
    \sum_{i=1}^N \|\Theta \bx^{(i)} - \by^{(i)}\|_2^2 
    =&  \sum_{i=1}^N \sum_{k=1}^m\|\theta_k^T \bx^{(i)} - \by^{(i)}_k\|_2^2\\
    =&  \sum_{k=1}^m\sum_{i=1}^N \|\theta_k^T \bx^{(i)} - \by^{(i)}_k\|_2^2.
\end{align}
We can solve the same linear regression problem for each $k$ that minimizes
\begin{align}
    \sum_{i=1}^N \|\theta_k^T \bx^{(i)} - \by^{(i)}_k\|_2^2.
\end{align}
then merge them to form $\Theta$.


%%% ERM
%%%%%%%%%%%%%%%%%%%%%%%%%%%%%%%%%%%%%%%%%%%%%%%%%%%%%%%%%%%%%%%%%%%%%%%
\section{Empirical Risk Minimization} \label{sec:ERM}
\subsection{Empirical Risk} \label{subsec:ER}
Given a data
\begin{align}
    (\bx^{(1)}, y^{(1)}), (\bx^{(2)}, y^{(2)}), \dots,  (\bx^{(N)}, y^{(N)})
\end{align}
that are sampled from $P$ in i.i.d.~manner where $\bx^{(i)}\in\cX$ and $y^{(i)}\in\cY$.
The classifier $f:\cX\rightarrow \cY$ is to predict $Y$ given $X$ for new data-label pair $(X, Y)\sim P$.


\begin{definition}
    A loss function $\ell:\fR^m\times \fR^m \rightarrow \fR$ quantifies how well $\hat{y}$ approximates $y$.
\end{definition}
Smaller $\ell(\hat{y}, y)$ indicates $\hat{y}$ is a good approximation of $y$.
Typlically, the loss function is nonnegative (but not always).

\begin{example}
    \begin{itemize}
        \item quadratic loss (Mean Squared Error)
        \begin{align}
            \ell(\hat{y}, y) = \|\hat{y} - y\|_2^2.
        \end{align}
        \item absolute loss ($\ell_1$-loss)
        \begin{align}
            \ell(\hat{y}, y) = \|\hat{y} - y\|_1.
        \end{align}
    \end{itemize}
\end{example}

Sometimes, we may want to predict $Y$ in different space $f:\cX\rightarrow \tilde{\cY}$.
One example is soft classification with cross-entropy loss.

\begin{definition}
    The expected risk (or true risk or population risk) is
    \begin{align}
        \cR[f] = \Expec{\ell(f(X), Y)}
    \end{align}
\end{definition}

Our goal is to solve 
\begin{align}
    \min_f \cR[f]
\end{align}

\begin{definition}
We call
\begin{align}
    \cR^\star = \inf_f \cR[f]
\end{align}
the Bayes risk or the optimal risk, where the infimum (minimum) is over all possible functions.
The optimal achieving $f^\star: \cX\rightarrow \tilde{\cY}$ is called Bayes predictor (Bayes optimal).
I.e.,
\begin{align}
    \cR^\star = \cR[f^\star].
\end{align}
\end{definition}

% \begin{theorem}
%     Let $f^\star$ be such that
%     \begin{align}
%         f^\star(x) \in \argmin _{\ty\in\tilde{\cY}} \Expec{\ell(\tilde{y}, Y)|X=x}
%     \end{align}
%     for all $x\in\cX$.
%     Then,
%     \begin{align}
%         \cR[f]\geq \cR[f^\star].
%     \end{align}
% \end{theorem}
% \begin{proof}
%    \begin{align}
%        \cR[f] =& \Expec{\Expec{\ell(f(X), Y)|X}}\\
%        \geq& \Expec{\Expec{\ell(f^\star(X), Y)|X}}\\
%        =&\cR[f^\star].
%    \end{align} 
% \end{proof}
% 
% \begin{example}[Binary Classification]
% Consider the binary classification problem where we have gray scale image $\bx^{(1)}, \dots, \bx^{(N)}\in\fR^d$,
% where $d$ is the number of pixels, and corresponding labels $y^{(1)}, \dots, y^{(N)}\in\{-1, +1\}$
% (E.g. labels indicate whether the image contains a cat or dog.)
% 
% Goal is to learn a function $f:\fR^d \rightarrow \{-1, +1\}$ that approximately solves
% \begin{align}
%     \min_f \Pr{f(\bX)\neq Y}.
% \end{align}
% In such case, we have a loss function 
% \begin{align}
%     \ell(f(\bX), Y) = \mathbf{1}\{f(\bX)\neq Y\}
% \end{align}
% so that 
% \begin{align}
%     \cR[f] =& \Expec{\ell(f(\bX), Y)}\\
%     =& \Expec{\mathbf{1}\{f(\bX)\neq Y}\}\\
%     =& \Pr{f(\bX)\neq Y}.
% \end{align}
% 
% The Bayes optimum is
%     \begin{align}
%         f^\star(x) =& \argmin _{\ty\in\{-1, +1\}} \Expec{\ell(\tilde{y}, Y)|X=x}\\
%         =& \argmin _{\ty\in\{-1, +1\}} \Pr{\tilde{y}\neq Y|X=x}\\
%         =& \argmax _{\ty\in\{-1, +1\}} \Pr{\tilde{y}= Y|X=x}\\
%         =& \argmax _{\ty\in\{-1, +1\}} P_{Y|X}(\tilde{y}|x).
%     \end{align}
% 
% The empirical risk is a fraction of error on training data.
% \end{example}
%     
% \begin{example}[Regression with squared loss]
%      
% \end{example}
% Note that both examples are basically the same as what we discussed in estimation.

Although our goal is to minimize the true risk $\cR[f]$, it is often impossible to compute since the true data distribution is unknown.
Instead, we use approximate 
\begin{align}
    \hat{\cR}[f] = \frac{1}{N}\sum_{i=1}^N \ell(f(\bx^{(i)}), y^{(i)})
\end{align}
We call $\hat{\cR}[f]$ the {\it\bf empirical risk}.
Note that $\hat{\cR}[f] \rightarrow \cR[f]$ as $N$ goes to infinity by the law of large numbers (LLN).

% \subsection{Measure of Performance} \label{subsec:measure of performance}
% We want to prove that our ML method (model) is good. 
% Although we often minimize the empirical risk, the true goal is to minimize the minimum expected risk (Bayes risk) $\cR^\star = \min_f \cR[f]$.
% Thus, the goal should be
% \begin{align}
%     \cR[f] - \cR^\star < \mbox{small}.
% \end{align}
% However, in many cases our $f$ is itself random , as it depends on randomly sampled data 
% and the randomness of the training algorithm (like SGD which we will discuss later).
% 
% Thus, the measure of performance can be expected error where the goal is
% \begin{align}
%     \Esub{f}{\cR[f] - \cR^\star} < \mbox{ small}.
% \end{align}
% This implies $f$ is good ``in expectation''.
% An alternative measure of performance is called {\it probably approximately correct (PAC)}:
% \begin{align}
%     \Pr{\cR[f] - \cR^\star < \mbox{ small}} > 1-\mbox{ small}.
% \end{align}
% The interpretation is that $f$ is good with high probability.
% PAC is a weaker notation (easier to prove) than expected error:
% \begin{itemize}
%     \item A bound on expected error implies PAC by Markov inequality
%     \item If an algorithm catastrophically fails with a small probability, then the expected error will be large while PAC holds.
% \end{itemize}
% 




% Many classifiers (or predictor) have the form $\hat{y} = f_\theta(x)$.
% The function class $\cF = \{f_\theta\}$ fixes the structure or form of the classifier.
% $\theta$ is a set of parameters, which can be a vector, matrix, or other structure.
% In linear regression problem, $\theta\in\fR^d$ is a vector and the model is
% \begin{align}
%     f_{\theta}(x) = \theta_1x_1 + \dots + \theta_d x_d.
% \end{align}
% 
% Choosing a particular $\theta$ given some training data
% \begin{align}
%     (\bx^{(1)}, y^{(1)}), (\bx^{(2)}, y^{(2)}), \dots,  (\bx^{(N)}, y^{(N)})
% \end{align}
% is called  ``training'' the model (or ``fitting'').
% We often assume that the data is sampled in i.i.d.~$\sim P$.


%%% Gradient Descent
%%%%%%%%%%%%%%%%%%%%%%%%%%%%%%%%%%%%%%%%%%%%%%%%%%%%%%%%%%%%%%%%%%%%%%%
\section{Gradient Descent} \label{sec:gradient}
This chapter heavily relies on~\cite{gower2018convergence, ryu2022large}.

%%% Convexity
\subsection{Convexity}
\begin{definition}
    A function $f:\fR^d\rightarrow \fR$ is convex if
    \begin{align}
        f(tx + (1-t)y) \leq tf(x) + (1-t) f(y)
    \end{align}
    for all $x, y\in\fR^d$ and $t\in[0, 1]$.
\end{definition}
If $\fR$ is differentiable and convex, then every tangent line lies below the function graph, i.e.,
\begin{align}
    f(y) \geq f(x) + \langle \nabla f(x) , y-x\rangle \label{eq:convex tangent}
\end{align}
for all $x, y\in\fR^d$, which can be directly derived from the definition.

\begin{remark}
    The above property is clear if $d=1$, where
    \begin{align}
        f(y) \geq f(x) + f'(x) (y-x),
    \end{align}
    which implies
    \begin{align}
        \frac{f(y)- f(x)}{y-x} \leq f'(x) & \mbox{, if $y< x$}\\
        \frac{f(y)- f(x)}{y-x} \geq f'(x) & \mbox{, if $y> x$}
    \end{align}
\end{remark}

Note that $\nabla f$ is a gradient of $f$, i.e.,
\begin{align}
    \nabla f(x) = \begin{bmatrix}
        \frac{\partial}{\partial x_1} f(x)\\
        \frac{\partial}{\partial x_2} f(x)\\
        \vdots \\
        \frac{\partial}{\partial x_d} f(x)
    \end{bmatrix}.
\end{align}
The gradient of $f$ is $d$-dimensional generalization of derivative.
A directional derivative measures the rate at which a function changes in a particular direction $v$ at a given point $x$,
\begin{align}
    \langle \nabla f (x) , v\rangle = \nabla f(x)^T v.
\end{align}


Furthermore, let $\nabla^2 f$ is a Hessian of $f$, where
\begin{align}
    \nabla^2 f(x) = \begin{bmatrix}
        \frac{\partial^2}{\partial x_1\partial x_1} f(x)& \frac{\partial^2}{\partial x_1\partial x_2} f(x)& \cdots & \frac{\partial^2}{\partial x_1\partial x_d} f(x)\\
        \frac{\partial^2}{\partial x_2\partial x_1} f(x)& \frac{\partial^2}{\partial x_2\partial x_2} f(x)& \cdots & \frac{\partial^2}{\partial x_2\partial x_d} f(x)\\
        & \vdots & & \\
        \frac{\partial^2}{\partial x_d\partial x_1} f(x)& \frac{\partial^2}{\partial x_d\partial x_2} f(x)& \cdots & \frac{\partial^2}{\partial x_d\partial x_d} f(x)
    \end{bmatrix}.
\end{align}
We have alternative definition for convex function $f$.
\begin{theorem}
        If $f:\fR^d \rightarrow \fR$ is twice-differentiable, the function $f$ is convex if and only if
        \begin{align}
            \nabla^2 f \succeq 0.
        \end{align}
        In other words, $\nabla^2f(x)$ is positive semi-definite for every $x$, and equivalently, $v^T \nabla^2 f(x) v \geq 0$ for all vector $v\in\fR^d$.
\end{theorem}


%%% Smootheness
\subsection{Smootheness}
\begin{definition}
    A function $f$ is $L$-Lipschitz continuous if
    \begin{align}
        |f(x_1) - f(x_2)| \leq L \cdot \|x_1-x_2\|
    \end{align}
    for all $x_1, x_2$.
\end{definition}
A Lipschitz continuous function is limited in how fast it can change.

\begin{definition}
    A differentiable function $f$ is said to be $L$-smooth if its gradients are $L$-Lipschitz continuous, i.e.,
    \begin{align}
        \|\nabla f(x_1) - \nabla f(x_2)\| \leq L \cdot \|x_1-x_2\|
    \end{align}
    for all $x_1, x_2$.
\end{definition}
An $L$-smooth function is limited in how fast its gradient can change.

\begin{theorem}
    If $f$ is twice-differentiable, then
    \begin{align}
        \nabla^2 f(x) \preceq LI
    \end{align}
    i.e., $LI - \nabla^2 f(x)$ is positive semidefinite matrix.
\end{theorem}
A rough proof sketch is following.
If $f$ is twice differentiable, then the Taylor expansion gives
\begin{align}
    \nabla f(x + tv) - \nabla f(x) = \int_{s=0}^t \nabla^2 f(x+sv)v\, ds
\end{align}
for all vector $v\in\fR^d$.
Since
\begin{align}
    \left\|\int_{s=0}^t \nabla^2 f(x+sv)v\, ds\right\|_2 \leq Lt \|v\|_2
\end{align}
and therefore,
\begin{align}
    \frac{\left\|\int_{s=0}^t \nabla^2 f(x+sv)v\, ds\right\|_2}{t \|v\|} \leq L.
\end{align}
Taking the limit as $t\rightarrow 0$, then
\begin{align}
    \nabla^2 f(x) \preceq L I.
\end{align}

\begin{exercise}
    Provide a rigorous proof of the above theorem.
\end{exercise}

The above theorem implies the following lemma.
\begin{lemma}\label{lem:lipschitz}
    For all $\delta\in \fR^d$
    \begin{align}
        f(x+\delta) \leq f(x) + \langle \nabla f(x), \delta\rangle + \frac{L}{2}\|\delta\|^2. 
    \end{align}
\end{lemma}

\begin{proof}
    Let $g(t) = f(x+ t(y-x))$.
    Then, $g'(t) = \nabla f(x+t(y-x))^T (y-x)$, and
    \begin{align}
        g'(t) - g'(0) =& \nabla f(x+t(y-x))^T (y-x) - \nabla f(x)^T(y-x)\\
        =&\frac{1}{t} (\nabla f(x+t(y-x)) - \nabla f(x))^T(x+t(y-x) - x)\\
        \leq & \frac{1}{t} L\|x+t(y-x) - x\|^2\\
        =& L t\|y-x\|^2.
    \end{align}
    Finally,
    \begin{align}
        f(y) =&g(1)\\
        =& g(0) + \int_0^1 g'(t)\, dt\\
        \leq&  g(0) + \int_0^1 (g'(0) + Lt \|y-x\|^2 )\, dt\\
        \leq&  g(0) + g'(0) + \frac{L}{2} \|y-x\|^2\\
        =& f(x) + \nabla f(x)^T (y-x) + \frac{L}{2} \|y-x\|^2.
    \end{align}
\end{proof}

\begin{remark}
    Note that the above theorem (and proof) does not require a convexity of function $f$, and it only requires the $L$-smootheness.
\end{remark}

\begin{theorem}
    If $f$ has a minimizer $x^\star$, where $f(x) \geq f(x^\star)$ for all $x$, then
    \begin{align}
        \frac{1}{2L} \|\nabla f(z)\|^2 \leq f(z) - f(x^\star) \leq \frac{L}{2} \|z-x^\star\|^2
    \end{align}
    for all $z\in\fR^d$.
\end{theorem}

\begin{proof}
    First, consider the right-hand inequality.
    \begin{align}
        f(z) \leq& f(x^\star) + \nabla f(x^\star)^T (z-x^\star) + \frac{L}{2}\|z-x^\star\|^2\\
        =& f(x^\star) + \frac{L}{2}\|z-x^\star\|^2
    \end{align}
    since $\nabla f(x^\star) = 0$ at the optimal point.

    Now consider the left-hand inequality.
    \begin{align}
        f(x^\star) =& \inf_y f(y)\\
        \leq& \inf_y f(z) + \nabla f(z)^T (y-z) + \frac{L}{2} \|y-z\|^2\\
        \leq& f(z) + \inf_{t, \|u\|=1} \nabla f(z)^T tu + \frac{L}{2} t^2\|u\|^2\\
        \leq& f(z) + \inf_{\|u\|=1} \inf_t (\nabla f(z)^T u) t + \frac{L}{2} t^2
    \end{align}
    where we re-parameterize $y = z + t u$ with a unit vector $u$.
    The infimum of the quadratic formula (with respect to $t$) is achieved by $t^\star = -\frac{1}{L}\nabla f(z)^T u$.
    Thus, 
    \begin{align}
        f(x^\star)
        \leq& f(z) + \inf_{\|u\|=1} (\nabla f(z)^T u) \left(-\frac{1}{L} \nabla f(z)^T u\right) + \frac{L}{2} \left(-\frac{1}{L} \nabla f(z)^T u\right)^2\\
        =& f(z) + \inf_{\|u\|=1} -\frac{1}{2L} \left(\nabla f(z)^T u\right)^2\\
        =& f(z) + \inf_{\|u\|=1} -\frac{1}{2L}  \|\nabla f(z)\|^2\\
        =& f(z) -\frac{1}{2L}  \|\nabla f(z)\|^2
    \end{align}
    which concludes the proof.
\end{proof}


%%% Smooth and Convex
\subsection{Smooth and Convex}
In many optimization problems, objective function is both smooth and convex.
\begin{lemma}
    If $f$ is convex and $L$-smooth, then
    \begin{align}
        f(y) - f(x) \leq& \langle \nabla f(y) , y-x \rangle - \frac{1}{2L}\|\nabla f(y) - \nabla f(x)\|^2\label{eq:coercivity1}\\
        \langle \nabla f(y)- \nabla f(x) , y-x \rangle \geq& \frac{1}{L}\| \nabla f(x) - \nabla f(y)\|^2\label{eq:coercivity2}
    \end{align}
    where the last inequlaity is often called Co-coercivity.
\end{lemma}

Co-coercivity means the gradient of a smooth convex function cannot twist or oscillate:
its change is aligned with the displacement,
making gradient descent (GD) stable, predictable, and convergent.

\begin{proof}

First, from the $L$-smoothness
\begin{align}
    f(x+z) \leq f(x) + \langle \nabla f(x), z\rangle + \frac{L}{2}\|z\|^2
\end{align}
and from convexity
\begin{align}
    f(y) + \langle \nabla f(y), z+x-y\rangle \leq f(x+z).
\end{align}
By adding the above two and rearrange terms, we have
\begin{align}
-\frac{L}{2}\|z\|^2 + \langle \nabla f(y) - \nabla f(x), z\rangle \leq f(x)- f(y) - \langle \nabla f(y), x-y\rangle.
\end{align}
By plugging $z = \frac{1}{L}(\nabla f(y) - \nabla f(x))$, we achieve the correct form of~\eqref{eq:coercivity1}.


    Let $g_x(z) = f(z) - \nabla f(x)^T z$, which is a function of $z$ (and you can think $x$ is constant).
    Since $f$ is a convex function $\nabla^2 g_x(z) \geq 0$ and $\nabla g_x(x) = 0$.
    Thus, $g_x$ is minimized at $z^\star=x$.
    Also, it is clear that $\nabla_z g_x(z) = \nabla_z(f(z) - \nabla f(x)^T z) = \nabla_z f(z) - \nabla f(x)$.
    Then, we have
    \begin{align}
        g_x(y) - g_x(x) \geq& \frac{1}{2L} \|\nabla_z g_x(y)\|^2\\
        =& \frac{1}{2L} \|\nabla f(y) - \nabla f(x)\|^2
    \end{align}
    If we similarly define $g_y(z) = f(z) - \nabla f(y)^T z$, then
    \begin{align}
        g_y(x) - g_y(y) \geq& \frac{1}{2L} \|\nabla f(x) - \nabla f(y)\|^2.
    \end{align}
    By adding two inequality, we have
    \begin{align}
        (g_x(y) - g_x(x)) + (g_y(x) - g_y(y))
        =& \left(f(y) -\nabla f(x)^T y  - f(x) + \nabla f(x)^Tx \right)\nonumber\\
        &+ \left(f(x) - \nabla f(y)^T x - f(y) +\nabla f(y)^T y\right)\\
        =& (\nabla f(x) - \nabla f(y))^T (x-y),
    \end{align}
    and
    \begin{align}
        (g_x(y) - g_x(x)) + (g_y(x) - g_y(y))
        \geq &\frac{1}{2L} \|\nabla f(y) - \nabla f(x)\|^2+ \frac{1}{2L} \|\nabla f(x) - \nabla f(y)\|^2\\
        \geq &\frac{1}{L} \|\nabla f(x) - \nabla f(y)\|^2.
    \end{align}
    Finally, 
    \begin{align}
        (\nabla f(x) - \nabla f(y))^T (x-y) \geq &\frac{1}{L} \|\nabla f(x) - \nabla f(y)\|^2.
    \end{align}
\end{proof}

%%% Strong Convexity
\subsection{Strong Convexity}
\begin{definition}
    A convex function $f$ is called $\mu$-strong convex if
    \begin{align}
        f(y)\geq f(x) + \langle \nabla f(x), y-x\rangle + \frac{\mu}{2} \|y-x\|^2
    \end{align}
    for all $x, y\in\fR^d$.
\end{definition}
It implies the function is approximately behave like quadratic function near $x$.

\begin{theorem}
    If a convex function $f$ is $\mu$-strongly convex, then
    \begin{align}
        \nabla^2 f(x) \succeq \mu I.
    \end{align}
\end{theorem}

\begin{theorem}
    If $f$ is differentiable and convex, then
    \begin{align}
        (\nabla f(x) - \nabla f(y))^T (x-y) \geq 0
    \end{align}
    for all $x, y$.
\end{theorem}
\begin{proof}
    We have
    \begin{align}
        f(x) \geq& f(y) + \nabla f(y) (x-y)\\
        f(y) \geq& f(x) + \nabla f(x) (y-x).
    \end{align}
    By adding two inequalities, we have
    \begin{align}
        (\nabla f(x) - \nabla f(y))^T (x-y) \geq 0.
    \end{align}
\end{proof}

\begin{theorem}
    If $f$ is $\mu$-strongly convex, 
    \begin{align}
        (\nabla f(x) - \nabla f(y))^T (x-y) \geq \mu \|x-y\|^2
    \end{align}
\end{theorem}

\begin{proof}
    We have
    \begin{align}
        f(x) \geq& f(y) + \nabla f(y) (x-y) + \frac{\mu}{2} \|x-y \|^2\\
        f(y) \geq& f(x) + \nabla f(x) (y-x) + \frac{\mu}{2} \|y-x \|^2.
    \end{align}
    By adding two inequalities, we have
    \begin{align}
        (\nabla f(x) - \nabla f(y))^T (x-y) \geq \mu\|x-y\|^2.
    \end{align}
\end{proof}



\begin{lemma}
    If $f$ is $\mu$-strongly convex, then it satisfies the Polyak-Lojasiewicz (PL) condition, which is
    \begin{align}
        \|\nabla f(x) \|^2 \geq 2\mu (f(x) - f(x^\star)).
    \end{align}
\end{lemma}
\begin{proof}
    \begin{align}
        f(x) - f(x^\star)\leq& \langle \nabla f(x), x- x^\star\rangle -\frac{\mu}{2} \|x^\star - x\|^2\\
        =& -\frac{1}{2} \|\sqrt{\mu}(x-x^\star) - \frac{1}{\sqrt{\mu}} \nabla f(x)\|^2 + \frac{1}{2\mu} \|\nabla f(x)\|^2\\
        \leq & \frac{1}{2\mu}\|\nabla f(x)\|^2.
    \end{align}
\end{proof}
Note that strong convexity is sufficient condition for Polyak-Lojasiewicz (PL) condition.
PL condition is an interesting codition by itself, which will guarantee the linear convergence of GD.



%%% Gradient Descent
\subsection{Gradient Descent}


Consider a differentiable function $f:\fR^d \rightarrow \fR$, and we want to solve
\begin{align}
    \mbox{minimize}_{x\in\fR^d} f(x).
\end{align}
Without any assumption, we are not able to find a global minimum $x^\star$.
However, if $f$ is convex, then $x$ is solution if and only if $\nabla f(x) = 0$.
A gradient descent algorithm is
\begin{align}
    x^{(k+1)} = x^{(k)} - \alpha_k \nabla f (x^{(k)})
\end{align}
where $\alpha_k>0$ is learning rate at $k$-th iteration.
The intuition of the algorithm is to update $x$ toward the maximally decreasing direction (since $\nabla f(x)$ is the maximally increasing direction of $f(x)$).

\begin{remark}
    If $d=1$, the motivation is clear.
    We are looking for a point $f'(x)=0$.
    Starting from a random point $x_0$, compute $f'(x_0)$.
    If $f'(x_0) > 0$, the function is increasing at $x=x_0$, and therefore search the region $x < x_0$ where the next iteration would be $x_1 = x_0 - \alpha f'(x_0)$.
    If $f'(x_0) < 0$, the function is decreasing at $x=x_0$, and therefore search the region $x > x_0$ where the next iteration would be $x_1 = x_0 - \alpha f'(x_0)$ (here we have $f'(x_0)<0$.
    The learning rate is a ``step-size'', how far your next iteration would be.
\end{remark}


We further assume that $\nabla f$ is $L$-Lipschitz continuous.
Then, the following theorem guarantees that the Gradient descent always converges with appropriate learning rate.
\begin{theorem}
    If $f$ has global minimun $x^\star$ and $f$ is $L$-smooth, then Gradient Descent algorithm converges to the local minimum (where $\nabla f(x) = 0$).
\end{theorem}

\begin{proof}
Then,
\begin{align}
    f(x^{(k+1)}) =& f(x^{(k)} - \alpha_k \nabla f(x^{(k)}))\\
    \leq& f(x^{(k)}) + \langle \nabla f(x^{(k)}), -\alpha_k \nabla f(x^{(k)})\rangle + \frac{L}{2} \|\alpha_k \nabla f(x^{(k)})\|^2\\
    \leq& f(x^{(k)}) + \left(\frac{L}{2}\alpha_k^2 - \alpha_k\right) \|\nabla f(x^{(k)})\|^2
\end{align}
If we pick $\frac{L}{2}\alpha_k ^2 - \alpha_k < 0 $ (e.g., $\alpha_k<2/L$), then
\begin{align}
    f(x^{(k+1)}) \leq& f(x^{(k)})
\end{align}
Since $f(x^{(k)})$ is lower bounded by $f(x^\star)$ (where $x^\star$ is a global minimum) and decreasing,
thus it converges to ``something''.
If $f(x^{(k)})$ converges to ``something,'' then $\|\nabla f(x^{(k)})\|^2$ should converge to 0.
\end{proof}


\begin{remark}
Do you think this theoretical result is too weak? or seems useless? 
(since we do not know whether the function is Lipschitz or not)
Indeed this theorem tells you the most important fact.
If your function (and its gradient) varies too quickly, you need to be careful (adjust step size).
Intuitively, if your function is varying crazily, no algorithm can guarantee the convergence.
This beautiful theorem basically shows the relation between steepness 
of the function and the right step size of the algorithm.
\end{remark}

% \begin{exercise}
%     Prove the Lemma~\ref{lem:lipschitz}.
% \end{exercise}
% 
% \begin{proof}
%     For any $y$ and $x$, from Taylor expansion,
%     \begin{align}
%         f(y) =& f(x) + \langle \nabla f(x) , y-x\rangle + \frac{1}{2} \nabla^2 f(\tx) \|y-x\|^2\\
%         \leq& f(x) + \langle \nabla f(x) , y-x\rangle + \frac{L}{2} \|y-x\|^2.
%     \end{align}
%     The last inequality is due to $L$-Lipschitzness of $\nabla f$.
% \end{proof}

\begin{exercise}
    Do you think the smoothness condition is necessary? Can we get rid of it?
    Also, can you interpret the meaning of $\alpha_k < 2/L$.
\end{exercise}

Consider the basic gradient step when $f$ is a convex function.
\begin{align}
    x^+ = x - \alpha \nabla f(x).
\end{align}
Then,
\begin{align}
    f(x^+) &= f(x-\alpha \nabla f(x)) \\
    &\leq f(x) +\nabla f(x)^T (x-\alpha \nabla f(x) - x) + \frac{L}{2} \|x-\alpha \nabla f(x) - x\|^2\\
    &= f(x) -\alpha (1- \alpha L/2) \|\nabla f(x)\|^2.
\end{align}
Let $0<\alpha \leq 1/L$, then
\begin{align}
    f(x^+) &\leq f(x) -\frac{\alpha}{2}  \|\nabla f(x)\|^2.
\end{align}
From the above inequality, we can say that the single gradient step decreases a function value since $f(x^+)\leq f(x)$.

From convexity,
\begin{align}
    \nabla f(x)^T (x^\star-x) \leq f(x^\star)-f(x).
\end{align}
By adding two inequalities,
\begin{align}
    f(x^+) - f(x^\star) \leq& \nabla f(x)^T (x- x^\star) -\frac{\alpha}{2} \|\nabla f(x)\|^2\\
    =&\frac{1}{2\alpha} \left(\|x-x^\star\|^2 - \|x-x^\star - \alpha \nabla f(x)\|^2\right)\\
    =&\frac{1}{2\alpha} \left(\|x-x^\star\|^2 - \|x^+-x^\star\|^2\right).
\end{align}
Since $f(x^\star) \leq f(x^+)$, we have
\begin{align}
    \|x-x^\star\|^2 \geq \|x^+-x^\star\|^2.
\end{align}
In other words, the distance to the optimal point also decreases in the gradient method.

If we iteratively apply the gradient method with initial point $x_0$ and constant step size $\alpha$, we have
\begin{align}
    \sum_{i=1}^T f(x_i) - f(x^\star) \leq& \frac{1}{2\alpha} \sum_{i=1}^T \left(\|x_{i-1}-x^\star\|^2 - \|x_i-x^\star\|^2\right)\\
    =& \frac{1}{2\alpha} (\|x_0-x^\star\|^2 - \|x_T - x^\star\|^2).\\
    \leq& \frac{1}{2\alpha} \|x_0-x^\star\|^2.
\end{align}
Since $f(x_i)$ is non-increasing (i.e., $f(x_i)\leq f(x_{i-1})$), we have
\begin{align}
    f(x_T)-f(x^\star) \leq& \frac{1}{T} \sum_{i=1}^T f(x_i) - f(x^\star) \\
    \leq& \frac{1}{2\alpha T} \|x_0 - x^\star\|^2.
\end{align}
Thus, the function value decreases with a rate of $O(1/T)$.

\begin{corollary}
    If you want to achieve $f(x_T) - f(x^\star)\leq \epsilon$, then the number of iteration should be $O(1/\epsilon)$.
\end{corollary}


%%% Gradient Descent
\subsection{Gradient Descent: $\mu$-strongly Convex Function}
Let $f$ be $\mu$-strongly convex and $L$-smooth.
\begin{theorem}
    Let $h(x) = f(x) - \frac{\mu}{2} \|x\|^2$, then $h$ is $L-\mu$-smooth.
\end{theorem}

\begin{proof}
Since
\begin{align}
    \nabla^2 f = \nabla^2 h + \mu I,
\end{align}
we have
\begin{align}
    \nabla^2 h = \nabla^2 f - \mu I \preceq (L-\mu) I.
\end{align}
\end{proof}

Since $h$ convex and $\nabla h (x) = \nabla f(x) - \mu x$
\begin{align}
    0\leq& (\nabla h(x) -\nabla h(y))^T (x-y)\\
    =& (\nabla f(x) - \nabla f(y))^T(x-y) - \mu\|x-y\|^2\\
    \leq& (L-\mu)\|x-y\|^2.
\end{align}
From co-coercivity of $h$, we also have
\begin{align}
    (\nabla h(x) - \nabla h(y))^T (x-y) \geq \frac{1}{L-\mu}\| \nabla h(x) - \nabla h(y)\|^2.
\end{align}
Equivalently,
\begin{align}
    &(\nabla f(x) - \nabla f(y) -\mu x + \mu y)^T (x-y) \geq \frac{1}{L-\mu}\| \nabla f(x) - \nabla f(y) - \mu x + \mu y\|^2\\
    &\Leftrightarrow 
    (\nabla f(x) - \nabla f(y))^T(x-y) -\mu \|x-y\|^2 \nonumber\\
    &\quad\geq \frac{1}{L-\mu}\| \nabla f(x) - \nabla f(y)\|^2 +\frac{\mu^2}{L-\mu} \|x-y\|^2- 2\frac{\mu}{L-\mu} (\nabla f(x) - \nabla f(y))^T(x-y)\\
    &\Leftrightarrow 
    \frac{L+\mu}{L-\mu} (\nabla f(x) - \nabla f(y))^T(x-y) \geq \frac{1}{L-\mu}\| \nabla f(x) - \nabla f(y)\|^2 +\frac{L\mu}{L-\mu} \|x-y\|^2\\
    &\Leftrightarrow 
    (\nabla f(x) - \nabla f(y))^T(x-y) \geq \frac{1}{L+\mu}\| \nabla f(x) - \nabla f(y)\|^2 +\frac{L\mu}{L+\mu} \|x-y\|^2.
\end{align}

Consider the gradient desecent method with $\mu$-strongly convex assumption on $f$, where
\begin{align}
    x^+ = x - \alpha\nabla f(x).
\end{align}
Let the step size $\alpha \leq \frac{2}{\mu+L}$, then from the above co-coercivity:
\begin{align}
    \|x^+-x^\star\|^2 =& \|x - \alpha \nabla f(x) - x^\star\|^2\\
    =& \|x- x^\star\|^2 \textcolor{red}{-2 \alpha \nabla f(x)^T (x-x^\star)} + \alpha^2 \|\nabla f(x)\|^2\\
    \leq& \|x- x^\star\|^2  \textcolor{red}{-2\alpha\frac{1}{L+\mu}\|\nabla f(x)\|^2 - 2\alpha\frac{L\mu}{L+\mu}\|x-x^\star\|^2} + \alpha^2 \|\nabla f(x)\|^2\\
    \leq& \left(1 - 2\alpha \frac{L\mu}{L+\mu} \right) \|x- x^\star\|^2  + \left(\alpha^2  - \frac{2\alpha}{L+\mu}\right)\|\nabla f(x)\|^2\\
    \leq& \left(1 - 2\alpha \frac{L\mu}{L+\mu} \right) \|x- x^\star\|^2.
\end{align}
which is because $\nabla f(x^\star) = 0$.

Let $c = 1 - 2\alpha \frac{L\mu}{L+\mu}<1$.
Then, the gradient descent is a ``contractive'' operation
\begin{align}
    \|x_i-x^\star\|^2 \leq c \|x_{i-1} - x^\star\|^2 \leq c^i \|x_0-x^\star\|^2.
\end{align}
Thus, now $x_i$ converges to $x^\star$ exponentially fast.
Furthermore
\begin{align}
    f(x_i) - f(x^\star) \leq \frac{L}{2}\|x_i- x^\star\|^2 \leq c^i \frac{L}{2} \|x_0 - x^\star\|^2.
\end{align}
Thus, the function value also converges to $f(x^\star)$ exponentially fast.

\begin{remark}
    If $f$ is $\mu$-strongly convex, the number of iterations to reach $f(x_i)-f(x^\star)\leq \epsilon$ is $O(\log 1/\epsilon)$.
\end{remark}


%%% Stochastic Gradient Descent
%%%%%%%%%%%%%%%%%%%%%%%%%%%%%%%%%%%%%%%%%%%%%%%%%%
\section{Stochastic Gradient Descent}
In many deep learning applications, the loss function is of form
\begin{align}
    \cL(\theta) = \frac{1}{N}\sum_{i=1}^N \cL_i(\theta) \label{eq:sgd objective}
\end{align}
For example, in linear regression problem,
we have $N$ data points $\{(x^{(1)}, y^{(1)}), \dots , (x^{(N)}, y^{(N)})\}$ 
with model $h_\theta$ and loss function $\ell(\cdot, \cdot)$.
The ERM is to minimize the following risk:
\begin{align}
    \cL(\theta) =    \frac{1}{N} \sum_{i=1}^N \ell(h_{\theta}(x^{(i)}), y^{(i)}),
\end{align}
where $h_\theta$ is a model, which is of form of \eqref{eq:sgd objective}.
In this section, we assume that the $N$ is extremely large 
(number of data is huge, for instance, LLM training).
Since computing the gradient of $N$ functions are challenging,
we often apply stochastic gradient descent (SGD) in such case.

In full gradient descent, one should compute
\begin{align}
    \theta_{k+1} =
    \theta_k - \alpha \frac{1}{N}\sum_{i=1}^N \nabla \ell(h_\theta(x^{(i)}), y^{(i)})
\end{align}
at $k$-th iteration.
The idea of stochastic gradient descent (SGD) is to approximate the true gradient
with a subsamples, i.e.,
\begin{align}
    \theta_{k+1} = \theta_k - \alpha  \nabla \ell(h_\theta(x^{(i_k)}), y^{(i_k)})
\end{align}
where $i_k$ is randomly chosen index from $\{1, 2, \dots, N\}$.
Note that 
\begin{align}
  \Expec{  \nabla \ell(h_\theta(x^{(i_k)}), y^{(i_k)})} = 
  \frac{1}{N}\sum_{i=1}^N \nabla \ell(h_\theta(x^{(i)}), y^{(i)})
\end{align}
and the estimate is unbiased.
In practice, the indices are usually chosen without replacement until we complete one full
cycle through the entire dataset.


%%% Mini-batching
%%%%%%%%%%%%%%%%%%%%%%%%%%%%%%%%%%%%%%%%%%%%%%%%%
\subsection{Mini-batching}
A common technique employed with SGD is mini-batching, where we choose a random subset $I_k\in\{1, \dots, N\}$ with size $|I_k|=b\ll N$.
We then repeat
\begin{align}
    \theta_{k+1} =
    \theta_k - \alpha \frac{1}{b}\sum_{i\in I_k} \nabla \ell(h_\theta(x^{(i)}), y^{(i)})
\end{align}
Since 
\begin{align}
  \Expec{  \frac{1}{b}\sum_{i\in I_k} \nabla \ell(h_\theta(x^{(i)}), y^{(i)})} = 
  \frac{1}{N}\sum_{i=1}^N \nabla \ell(h_\theta(x^{(i)}), y^{(i)})
\end{align}
we still have an unbiased estimate of the full gradient.
Furthermore, we not reduced the variance of our gradient estimate by $1/b$,
but we have also incurred $b$ times more computational cost.
In practice, the batch-update can be parallelized and we then indeed gain computational savings.

%%% Analysis
%%%%%%%%%%%%%%%%%%%%%%%%%%%%%%%%%%%%%%%%%%%%%%%%%
\subsection{Analysis}
Now let 
Let $g_k$ be a stochastic gradient of $\cL$ at $\theta_k$.
We assume that it is unbiased
\begin{align}
    \Esub{k}{g_k} = \Expec{g_k|\theta_0, \dots, \theta_k} = \nabla \cL(\theta_k)
\end{align}
where the expectation is conditioned on the iterates up to $\theta_k$.
We also assume that the convditional variance is bounded:
\begin{align}
    \Esub{k}{(g_k-\Esub{k}{g_k})^2} \leq \sigma^2.
\end{align}
This implies that the stochastic gradient is not too noisy (or close to the true gradient).

\begin{theorem}
    Let $\cL(\theta)$ be $L$-Lipschitz continuous convex function where the minimzer $\theta^\star$ exists.
    Let $\theta_0$ be a starting point and $K>0$ be the total iteration count.
    Then, Stochastic Gradient Descent (SGD) with the constant learning rate
    \begin{align}
        \alpha_k = \alpha = \frac{\|\theta_0 - \theta^\star\|}{\sqrt{L^2+\sigma^2}\sqrt{K+1}}
    \end{align}
    exhibits the rate
    \begin{align}
        \Expec{\cL(\bar{\theta}_K) - \cL(\theta^\star)} \leq \frac{\sqrt{L^2+\sigma^2}\|\theta_0-\theta^\star\|}{\sqrt{K+1}}
    \end{align}
    where
    \begin{align}
        \bar{\theta}_K = \frac{1}{K+1} \sum_{k=0}^K \theta_k.
    \end{align}
\end{theorem}

\begin{proof}
    First,
    \begin{align}
        \Esub{k}{\|\theta_{k+1}-\theta^\star\|^2} 
        =& \|\theta_k - \theta^\star\|^2 - 2\alpha \langle \Esub{k}{g_k}, \theta_k-\theta^\star\rangle + \alpha^2 \Esub{k}{\|g_k\|^2}\\
        \leq& \|\theta_k - \theta^\star\|^2 - 2\alpha (\cL(\theta_k) - \cL(\theta^\star)) + \alpha^2 (L^2+\sigma^2).
    \end{align}
    We take the total expectation on bith sides to get
    \begin{align}
        \Expec{\|\theta_{k+1}-\theta^\star\|^2} 
        \leq& \Expec{\|\theta_k - \theta^\star\|^2} - 2\alpha \Expec{\cL(\theta_k) - \cL(\theta^\star)} + \alpha^2 (L^2+\sigma^2).
    \end{align}
    By summing up from $k=0$ to $K$, we have
    \begin{align}
        \Expec{\|\theta_{K}-\theta^\star\|^2} 
        \leq& \Expec{\|\theta_0 - \theta^\star\|^2} - \sum_{k=0}^{K} 2\alpha \Expec{\cL(\theta_k) - \cL(\theta^\star)} + (K+1)\alpha^2 (L^2+\sigma^2).
    \end{align}
    I.e.,
    \begin{align}
    \frac{1}{K+1} \sum_{k=0}^{K}  \Expec{\cL(\theta_k) - \cL(\theta^\star)} 
\leq& \frac{1}{(K+1)2\alpha}\Expec{\|\theta_0 - \theta^\star\|^2} -\frac{1}{(K+1)2\alpha}\Expec{\|\theta_{K+1}-\theta^\star\|^2} 
+ \frac{\alpha}{2} (L^2+\sigma^2)\\
\leq& \frac{1}{(K+1)2\alpha}\|\theta_0 - \theta^\star\|^2 + \frac{\alpha}{2} (L^2+\sigma^2)\\
=& \frac{1}{\sqrt{K+1}\sqrt{L^2+\sigma^2}}\|\theta_0 - \theta^\star\|
    \end{align}
    On the other hand, due to the convexity, we have
    \begin{align}
    \frac{1}{K+1} \sum_{k=0}^{K}  \Expec{\cL(\theta_k) - \cL(\theta^\star)} 
    \geq 
    \Expec{\cL(\bar{\theta}_K) - \cL(\theta^\star)}.
    \end{align}
    Finally, we have
    \begin{align}
    \Expec{\cL(\bar{\theta}_K) - \cL(\theta^\star)}
\leq& \frac{1}{\sqrt{K+1}\sqrt{L^2+\sigma^2}}\|\theta_0 - \theta^\star\|
    \end{align}
\end{proof}

\begin{remark}
The convergence rate of Stochastic Gradient Descent of convex function is $O(1/\sqrt{K})$.
\end{remark}

\begin{remark}
In stochastic gradient descent (SGD), the theoretical proof focuses on $\bar{\theta}_K$ rather than $\theta_K$.
This means SGD is applied over $K$ iterations, and the average of $\theta_k$ values from $k = 0$ to $K$ is considered. 
However, in practice, it is common to use only the final iterate, $\theta_K$.
\end{remark}


%%% Regularized Method
%%%%%%%%%%%%%%%%%%%%%%%%%%%%%%%%%%%%%%%%%%%%%%%%%
\section{Regularized Method}
Consider the predictor (or classifier) $f_\theta(x)$ with parameter $\theta$.
An important attribute of predictor $f$ is ``sensitivity'' or ``continuity''.
We say $f_\theta$ is insensitive if for $x$ near $\tilde{x}$, the prediction $f_\theta(x)$ is near $f_\theta(\tilde{x})$.
Insensitive predictors often generalize well, especially when we do not have a lot of training data.
This is because there is a gap between the population risk
\begin{align}
    \Expec{\ell(f_\theta(X), Y)}
\end{align}
versus the empirical risk
\begin{align}
    \sum_{i=1}^N \ell(f_\theta(x^{(i)}), y^{(i)}).
\end{align}
More precisely, the predictor shows a good performance in empirical risk, 
but not in population risk.
Thus, the insensitive predictor is desirable.

A regularizer is a function $r:\fR^d \rightarrow \fR$ that measures the sensitivity of $f_\theta$.
In other words, if $r(\theta)$ is small then $f_\theta$ is insensitive.

\begin{example}
    In linear regression $f_\theta(x) = \theta^T x$, small sensitivity is associated with small $\theta$.
    By Cauchy-Schwartz inequality
    \begin{align}
        \|f_\theta(x) - f_\theta(\tx)\|_2 \leq\|\theta\|_2 \cdot \|x-\tilde{x}\|_2.
    \end{align}
    This suggest a regularizer $r(\theta) = \|\theta\|_2$
\end{example}

The common choice of regularizer is $\ell_2$ or square or ridge regularization:
\begin{align}
    r(\theta) = \sum_{i=1}^d \theta_i^2.
\end{align}
Another possible choice is the $\ell_1$-regularizer where
\begin{align}
    r(\theta) = \sum_{i=1}^d \theta_i^2.
\end{align}

With the above regularizer, we have the ridge empirical risk minmization (RERM) problem that minimizes
\begin{align}
    \sum_{i=1}^N \ell(f_\theta(x^{(i)}), y^{(i)}) + \lambda \|\theta\|^2.
\end{align}
As $\lambda$ increases, we have more insensitive solution.
    

\begin{exercise}
    The least square and ridge regression objective function is given by
    \begin{align}
        \|X\theta - Y\|^2 + \lambda \|\theta\|^2.
    \end{align}
    If $X$ is a full rank matrix, show that the minimum achiever is
    \begin{align}
        \theta^\star = (X^TX+\lambda I)^{-1}X^T Y.
    \end{align}
\end{exercise}

%%%%%%%%%%%%%%%%%%%%%%%%%%%%%%%%%%%%%%%%%%%%%%%%%
\subsection{Strongly convex SGD}
Consider the stochastic optimization problem
\begin{align}
    \min f(x) + \frac{\mu}{2} \|x\|^2.
\end{align}
Assume that $F$ is Lipschitz continuous.
Consider the stochastic gradient descent (SGD)
\begin{align}
    x_{k+1} = x_k - \alpha_k (g_k + \mu x_k)
\end{align}
where $g_k$ is stochastic gradient of $f$ at $x_k$.

\begin{theorem}
    Let $f:\fR^d\rightarrow \fR$ be a $G$-Lipschitz continuous convex function and $x^\star$ be a miminizer of $f(x) + \mu/2 \|x\|^2$.
    Suppose $\Esub{k}{g_k} = \nabla f(x_k)$ and $\mbox{Var}_k(g_k)\leq \sigma^2$, then SGD with step size
    \begin{align}
        \alpha_k = \frac{1}{\mu (k+1)}
    \end{align}
    exhibits the rate
    \begin{align}
        f(\bar{x}_k)-f(x^\star) \leq \frac{2(G^2+\sigma^2)}{\mu} \frac{1+\log (k+1)}{k+1}
    \end{align}
    where 
    \begin{align}
        \bar{x}_k = \frac{1}{k+1}\sum_{i=0}^k x_i.
    \end{align}
\end{theorem}

For simplicity, let assume $x_0=0$.
\begin{align}
    x_{k+1} =& x_k - \alpha_k (g_k + \mu x_k)\\
    =& (1-\alpha_k \mu)x_k + \alpha_k \mu (-\frac{1}{\mu}g_k)\\
    =& \sum_{i=0}^k \theta_i (-\frac{1}{\mu} g_i)
\end{align}
where $\theta_i = \alpha_i \mu\prod_{s=i+1}^k(1-\alpha_s \mu) $ and $\sum_{i=0}^k \theta_i = 1$.
Using Jensen's inequality,
\begin{align}
    \Expec{\|g_k + \mu x_k\|^2}
    &\leq 2\Expec{\|g_k\|^2} + 2 \mu^2 \Expec{\|x_k\|^2}\\
    &\leq 2(G^2+\sigma^2) + 2 \mu^2 \sum_{i=0}^k \theta_i \Expec{\|(-1/\mu)g_i\|^2}\\
    &= 2(G^2+\sigma^2) + 2(G^2+\sigma^2) = 4(G^2+\sigma^2).
\end{align}
Finally,
\begin{align}
    \Esub{k}{\|x_{k+1}-x^\star\|^2}
    =& \|x_k-x^\star\|^2 - 2\alpha_k\langle \Esub{k}{g_k} + \mu x_k, x_k-x^\star\rangle + \alpha_k^2 \Esub{k}{\|g_k+\mu x_k\|^2}\\
    \leq & \|x_k-x^\star\|^2 - 2\alpha_k (f(x_k) - f(x^\star) + \frac{\mu}{2} \|x_k-x^\star\|^2) + \alpha_k^2 4(G^2+\sigma^2)\\
    \leq & (1-\alpha_k \mu) \|x_k-x^\star\|^2 - 2\alpha_k (f(x_k) - f(x^\star)) + \alpha_k^2 4(G^2+\sigma^2).
\end{align}
Rearranging the terms and plugging in $\alpha_k = \frac{1}{\mu(k+1)}$, 
\begin{align}
    f(x_k)- f(x^\star) \leq \frac{\mu k}{2} \|x_k-x^\star\|^2 - \frac{\mu (k+1)}{2} \Esub{k}{\|x_{k+1}-x^\star\|^2} 
    + \frac{2(G^2+\sigma^2}{\mu(k+1)}.
\end{align}
Then, from the telescoping sum argument, we have
\begin{align}
    \sum_{i=0}^k f(x_i)- f(x^\star) 
    \leq&  \sum_{i=0}^k  \frac{2(G^2+\sigma^2)}{\mu(i+1)}\\
    \leq&  \frac{2(G^2+\sigma^2)}{\mu} (1+\log (k+1))\\
\end{align}
Finally,
\begin{align}
    f(\bar{x}_k) - f(x^\star) \leq \frac{1}{k} \sum_{i=0}^k f(x_i)- f(x^\star) 
    \leq  \frac{2(G^2+\sigma^2)}{\mu k} (1+\log (k+1))\\
\end{align}
The convergence rate is
\begin{align}
    \mathcal{O}\left(\frac{\log k}{k}\right).
\end{align}
This is faster than the SGD convergence rate 
\begin{align}
    \mathcal{O}\left(\frac{1}{\sqrt{k}}\right).
\end{align}

\begin{remark}
    The optimal rate (can be achieved and cannot do better than) is
\begin{align}
    \mathcal{O}\left(\frac{1}{k}\right).
\end{align}
The above SGD is slightly suboptimal with factor of $\log k$.
\end{remark}

%%% Gradient Descent as an Operator
\section{Operator Theory}
Let $\fT:\fR^d\rightarrow \fR^d$ be an operator.
\begin{definition}
    We say an operator $\fT$ is nonexpansive if
    \begin{align}
        \|\fT x - \fT y \| \leq \|x-y\|
    \end{align}
    for all $x, y$.
\end{definition}
Note that $\fT$ is nonexpansive if it is 1-Lipschitz.
Furthermore, if we say $\fT$ is a contraction if it is $L$-Lipschitz with $L< 1$.
\begin{definition}
    We say an operator $\fT$ is contractive if
    there exists $L<1$ such that 
    \begin{align}
        \|\fT x - \fT y \| \leq L \|x-y\|
    \end{align}
    for all $x, y$.
\end{definition}


\begin{definition}
For $\theta\in(0, 1)$, we say an operator $\fT$ is $\theta$-averaged if
    \begin{align}
        \fT = (1-\theta)\fI + \theta \fS
    \end{align}
    for an identity operator $\fI$ and some nonexpansive operator $\fS$.
\end{definition}

\begin{exercise}
    Show that $\theta$-averaged operator $\fT = (1-\theta)\fI + \theta\fS$ is also nonexpansive. 
\end{exercise}

\begin{definition}
    We say $x$ is a fixed point of $\fT$ if $x = \fT x$.
\end{definition}
The algorithm fixed-point iteration, also called the Picard iteration is
\begin{align}
    x^{(k+1)} = \fT x^{(k)}
\end{align}
for $k\geq 0$ with some starting point $x^{(0)}$.

\begin{remark}
    If $\fT$ is a contraction with Lipschitz constant $L<1$, then
    \begin{align}
        \|x^{(k)} - x^\star\| \leq L^k \|x^{(0)} - x^\star\|
    \end{align}
    for fixed point $x^\star$.
\end{remark}

\begin{theorem}
    Suppose $\fT$ is $\theta$-averaged with $\theta\in(0, 1)$, and there exists a fixed point of $\fT.$
    Then the fixed-point iteration converge with any starting point converges to some fixed point $x^\star$,
    and
    \begin{align}
        \|x^{(k+1)} - x^{(k)}\|^2 \leq& \frac{1}{k+1} \frac{\theta}{1-\theta} \|x^{(0)}- x^\star\|^2.
    \end{align}
\end{theorem}

\begin{proof}
    We have
    \begin{align}
        x^{(k+1)} = \fT x^{(k)} = (1-\theta) x^{(k)} + \theta \fS x^{(k)}.
    \end{align}
    Note that for all $x, y$, the following identity holds
    \begin{align}
        \|(1-\theta)x + \theta y\|^2 = (1-\theta) \|x\|^2 + \theta \|y\|^2 - \theta(1-\theta) \|x-y\|^2.
    \end{align}
    Then, for a fixed point $x^\star$, we have
    \begin{align}
        \|x^{(k+1)} - x^\star\|^2 
        =& (1-\theta) \|x^{(k)} - x^\star\|^2 + \theta \|\fS x^{(k)} - x^\star\|^2 - \theta (1-\theta) \|\fS x^{(k)} - x^{(k)}\|^2\\
        \leq & (1-\theta) \|x^{(k)} - x^\star\|^2 + \theta \|x^{(k)} - x^\star\|^2 - \theta (1-\theta) \|\fS x^{(k)} - x^{(k)}\|^2\\
        = &  \|x^{(k)} - x^\star\|^2 - \theta (1-\theta) \|\fS x^{(k)} - x^{(k)}\|^2. \label{eq:gd operator 1-step}
    \end{align}
    This implies the monotone decreases $\|x^{(k+1)}-x^\star\| \leq \|x^{(k)} - x^\star\|$.

    Consider the fixed-point residual $\fT x^{(k)} - x^{(k)}$.
    Since $\fT$ is nonexpansive,
    \begin{align}
        \|x^{(k+1)} - x^{(k)}\| \leq \|x^{(k)} - x^{(k-1)}\|.
    \end{align}
    By summing \eqref{eq:gd operator 1-step},
    \begin{align}
        \|x^{(k+1)} - x^{\star}\|^2 \leq \|x^{(0)} - x^\star\|^2 - \frac{1-\theta}{\theta} \sum_{j=0}^k \|\fT x^{(j)} - x^{(j)}\|^2.
    \end{align}
    This is equivalent to
    \begin{align}
        \sum_{j=0}^k \|\fT x^{(j)} - x^{(j)}\|^2 \leq \frac{\theta}{1-\theta} \|x^{(0)}- x^\star\|^2 - \frac{\theta}{1-\theta}\|x^{(k+1)}-x^\star\|^2.
    \end{align}
    Since $\|\fT x^{(j)} - x^{(j)}\|$ monotonically decreases, we have
    \begin{align}
        (k+1)\|x^{(k+1)} - x^{(k)}\|^2 \leq& \sum_{j=0}^k \|\fT x^{(j)} - x^{(j)}\|^2\\
        \leq& \frac{\theta}{1-\theta} \|x^{(0)}- x^\star\|^2.
    \end{align}
\end{proof}

The gradient descent algorithm is
\begin{align}
    x^{(k+1)} = \fT x^{(k)}
\end{align}
where 
\begin{align}
    \fT = \fI - \alpha \nabla f
\end{align}
for a learning rate $\alpha$.

Consider an operator $\fS$
\begin{align}
    \fS = \fI - \frac{2}{L} \nabla f.
\end{align}
From $L$-smootheness of $f$,
\begin{align}
    \|\fS x - \fS y\|^2
    =& \|x - (2/L)\nabla f(x) - (y - (2/L)\nabla f(y))\|^2\\
    =& \|x-y\|^2 - \frac{4}{L} \left(\langle x-y, \nabla f(x) - \nabla f(y)\rangle -\frac{1}{L}\|\nabla f(x) - \nabla f(y)\|^2\right)\\
    \leq& \|x-y\|^2.
\end{align}
Thus, $\fS$ is nonexpansive, and $\fT$ is $\theta$-averaged operator for $\alpha\in(0, 2/L)$ with $\theta=\alpha L/2 < 1$.
Consequently $x^{(k)}$ converges to a fixed point $x^\star$, where the fixed point $x^\star$ satisfies
\begin{align}
    \nabla f(x^\star) = 0.
\end{align}
If there exists a fixed point, then
\begin{align}
    \|\nabla f(x^{(k)})\|^2 = O(1/k).
\end{align}

Further, if $f$ is $\mu$-strongly convex and $\nabla f$ is $L$-Lipschitz, we can show the iteration is a contraction.
By the fundamental theorem of calculus,
\begin{align}
    \fT x - \fT y
    &= (x-y) - \alpha\bigl(\nabla f(x) - \nabla f(y)\bigr) \\
    &= (x-y) - \alpha \int_0^1 \nabla^2 f(tx + (1-t)y)(x-y)\,dt \\
    &= \int_0^1 \left[\fI - \alpha\nabla^2 f(tx+(1-t)y)\right](x-y) \,dt.
\end{align}
Then,
\begin{align}
    \|\fT x - \fT y\|
    &=\left\|\int_0^1 \left[\fI - \alpha\nabla^2 f(tx+(1-t)y)\right](x-y) \,dt\right\|\\
    &\leq  \int_0^1 \left\|\left[\fI - \alpha\nabla^2 f(tx+(1-t)y)\right](x-y) \right\|\,dt \\
    &\le \int_0^1 \left\|\fI - \alpha\nabla^2 f(tx+(1-t)y)\right\|\,dt \,\|x-y\|.
\end{align}
Since $\mu I\preceq \nabla^2 f(z) \preceq L I$ for all $z$, every eigenvalue
$\lambda$ of $\nabla^2 f(z)$ lies in $[\mu,L]$, so the eigenvalues of
$\fI - \alpha \nabla^2 f(z)$ lie in $\{1-\alpha\lambda : \lambda\in[\mu,L]\}$.
Therefore
\begin{align}
    \left\|\fI - \alpha\nabla^2 f(z)\right\|
    \le \max_{\lambda\in[\mu,L]} |1-\alpha\lambda|
    = \max\{|1-\alpha\mu|, |1-\alpha L|\}.
\end{align}
Hence
\begin{align}
    \|\fT x - \fT y\|
    \leq \max\{|1-\alpha\mu|, |1-\alpha L|\}\,\|x-y\|.
\end{align}
In particular, if $0 < \alpha < 2/L$, then $\max\{|1-\alpha\mu|, |1-\alpha L|\} < 1$
and $\fT$ is a contraction.

%%% Acceleration Methods
%%%%%%%%%%%%%%%%%%%%%%%%%%%%%%%%%%%%%%%%%%%%%%%%%%
\section{Acceleration Methods}
For convex and $L$-smooth function $f$, the gradient descent method provides an $O(1/k)$ rate.
It is natural to ask whether we can increase the speed of convergence, and the answer is yes (the worst-case rate).
In optimization, an acceleration is a modification of a base method to improve the convergence rate.

Consider the problem
\begin{align}
    \min_{x\in\fR^n} f(x)
\end{align}
where $f$ is convex and $L$-smooth.
The method 
\begin{align}
    x_{k+1} =& y_k - \frac{1}{L}\nabla f(y_k)\\
    y_{k+1} =& x_{k+1}+\frac{k-1}{k+2}(x_{k+1}-x_k)
\end{align}
where $x_0=y_0$, is called Nesterov's accelerated gradient method (AGM).

\begin{lemma}\label{lem:equiv AGM}
    The following formula is equivalent to AGM.
\begin{align}
    x_{k+1} =& y_k - \frac{1}{L}\nabla f(y_k)\\
    z_{k+1} =& z_k - \frac{k+1}{2L}\nabla f(y_k)\\
    y_{k+1} =& \left(1-\frac{2}{k+2}\right) x_{k+1}+\frac{2}{k+2} z_{k+1}.
\end{align}
where $x_0=y_0 = z_0$.
\end{lemma}
\begin{exercise}
    Show that the above process in Lemma~\ref{lem:equiv AGM} provides the same $x_k$ and $y_k$ sequences.
\end{exercise}

\begin{theorem}
    Assume the convex, $L$-smoothe function $f$ has a global minimizer $x^\star$.
    Then, AGM converges with the rate
    \begin{align}
        f(x_k) - f(x^\star) \leq\frac {2L}{k^2} \|x_0-x^\star\|^2.
    \end{align}
\end{theorem}

\begin{proof}
    Let $\theta_k = \frac{k+1}{2}$, then it satisfies
    \begin{align}
        \theta_k^2 - \theta_k \leq \theta_{k-1}^2.
    \end{align}
    
    First, $f(x)$ is $L$-smooth implies $f(x) - \frac{L}{2}\|x-y_k\|^2$ is concave.
    This is because $\nabla^2 f - L I \preceq 0$.
    The concavity implies
    \begin{align}
        f(x) - \frac{L}{2}\|x-y_k\|^2 \leq f(y_k) + \nabla f(y_k)^T (x-y_k).
    \end{align}
    Instead of $x$, we plug in $x = x_{k+1} =- y_k - \frac{1}{L}\nabla f(y_k)$, then we have
    \begin{align}
        f(x_{k+1}) - f(y_k ) + \frac{1}{2L}\|\nabla f(y_k)\|^2\leq& 0. 
    \end{align}
    
    From convexity of $f$, we have
    \begin{align}
        f(y_k) - f(x_k)\leq& \nabla f(y_k)^T (y_k-x_k)\\
        f(y_k) - f(x^\star)\leq& \nabla f(y_k)^T (y_k-x^\star).
    \end{align}

    Note that 
    \begin{align}
        z_{k+1} = x_k - \frac{\theta_k}{L}\nabla f(y_k)
    \end{align}
    and 
    \begin{align}
        \|z_{k+1}-x^\star\|^2 =&  \| z_k - \frac{\theta_k}{L}\nabla f(y_k)-x^\star\|^2 \\
        =& \|z_k - x^\star\|^2 -2\frac{\theta_k}{L} \nabla f(y_k)^T (z_k - x^\star) + \frac{\theta_k^2}{L^2} \|\nabla f(y_k)\|^2.
    \end{align}
    Define
    \begin{align}
        V_k =& \theta_{k-1}^2 (f(x_k)-f(x^\star)) + \frac{L}{2}\|z_k-x^\star\|^2.
    \end{align}
    Then,
    \begin{align}
        V_{k+1}-V_k =&  \theta_{k}^2 (f(x_{k+1})-f(x^\star)) + \frac{L}{2}\|z_{k+1}-x^\star\|^2
        -\theta_{k-1}^2 (f(x_k)-f(x^\star)) - \frac{L}{2}\|z_k-x^\star\|^2\\
        =& \theta_k^2 \left( f(x_{k+1})- f(x^\star) + \frac{1}{2L}\|\nabla f(y_k)\|^2\right) -\theta_k \nabla f(y_k)^T (z_k - x^\star)  -\theta_{k-1}^2 (f(x_k)-f(x^\star)) \\
        \leq& \theta_k^2 \left( f(y_{k})- f(x^\star)\right) -\theta_k \nabla f(y_k)^T (z_k - x^\star)  -\theta_{k-1}^2 (f(x_k)-f(x^\star)) \\
        =& (\theta_k^2 -\theta_k)( f(y_k) - f(x_k)) + \theta_k (f(y_k)- f(x_k)) + (\theta_k^2-\theta_{k-1}^2)(f(x_k)-f(x^\star))\nonumber\\
        &-\theta_k \nabla f(y_k)^T (z_k-x^\star)\\
        \leq& (\theta_k^2 -\theta_k)( f(y_k) - f(x_k)) + \theta_k (f(y_k)- f(x^\star)) -\theta_k \nabla f(y_k)^T (z_k-x^\star)\\
        \leq& (\theta_k^2 -\theta_k)\nabla f(y_k)^T (y_k-x_k) + \theta_k \nabla f(y_k)^T(y_k-x^\star) - \theta_k \nabla f(y_k)^T (z_k-x^\star)\\
        \leq& \theta_k \nabla f(y_k)^T \left((1-\theta_k)x_k + \theta_k y_k - z_k\right)\\
        =& 0,
    \end{align}
    where the last equality is from the definition of $z_k$. 
    \begin{align}
        y_k = \left(1-\frac{1}{\theta_k}\right) x_k + \frac{1}{\theta_k}z_k.
    \end{align}
    Thus, $V_k$ is a non-increasing sequence
    \begin{align}
        V_k\leq V_{k-1}\leq\dots \leq V_0.
    \end{align}
    Since $\theta_{-1} = 0$, we have
    \begin{align}
        \theta_{k-1}^2 (f(x_k) - f(x^\star)) \leq V_k = f(x_k)-f(x^\star) + \frac{L}{2} \|z_k - x^\star\|^2 \leq \frac{L}{2} \|z_0-x^\star\|^2.
    \end{align}
    Since $z_0 = x_0$, we finally get
    \begin{align}
        f(x_k) - f(x^\star)  \leq \frac{2L}{k^2} \|x_0-x^\star\|^2.
    \end{align}
    
\end{proof}

The following is beyond the scope of our lecture but is good to know.
We call ``First-order method'' be any iteratie algorithm that selects $x_{k+1}$ in the set
\begin{align}
    x_0 + \mbox{span} \{ \nabla f(x_0), \nabla f(x_1), \dots, \nabla f(x_k)\}.
\end{align}
Consider the function class $f$ that satisfies the followings.
\begin{enumerate}
    \item $f$ is convex and differentiable
    \item $f$ is $L$-smooth
    \item The optimal $x^\star$ exists.
\end{enumerate}
Then, the following results shows the converse (impossibility).

\begin{theorem}
For every integer $k\leq\frac{d-1}{2}$ and every $x_0$, there existsw a function in the class that for any first-order method
\begin{align}
    f(x_k) - f(x^\star) \geq \frac{2}{32} \frac{L \|x_0-x^\star\|^2}{(k+1)^2}
\end{align}
\end{theorem}





% \subsection{Gradient Descent Momentum}
% Gradient Descent with Momentum is an optimization technique designed to accelerate convergence by incorporating information 
% about the previous steps in the optimization process.
% While standard gradient descent only considers the current gradient to update the parameters,
% momentum takes into account the accumulated gradient history to smooth out updates.
% This approach can be intuitively understood through the following analogy:
% 
% Imagine a ball rolling down a hill.
% In standard gradient descent, the ball moves in the direction of the steepest slope at each step,
% adjusting its position based solely on the current slope.
% While effective, this method can be slow and may oscillate significantly 
% when the slope alternates between steep and shallow directions (e.g., in valleys or ridges).
% 
% With momentum, the ball gains velocity as it rolls, combining the effects of previous gradients.
% This is akin to the ball being subject to both gravity (the current gradient) and inertia (its accumulated velocity).
% The velocity helps the ball overcome small bumps or oscillations, enabling it to travel faster towards the bottom of the hill.
% 
% The mathematical formulation of momentum reflects this idea.
% Let \(\mathbf{x}_t\) represent the position (parameters) at step \(t\). The updates are governed by:
% \[
% \mathbf{v}_{t+1} = \beta \mathbf{v}_t - \eta \nabla f(\mathbf{x}_t),
% \]
% \[
% \mathbf{x}_{t+1} = \mathbf{x}_t + \mathbf{v}_{t+1}.
% \]
% Here:
% \begin{itemize}
%     \item \(\mathbf{v}_t\) is the velocity term, which combines the accumulated gradient history.
%     \item \(\beta \in [0, 1)\) is the momentum coefficient, determining how much past gradients influence the current update.
%     \item \(\eta\) is the learning rate, controlling the step size.
% \end{itemize}
% 
% Benefits of Gradient Descent Momentum is the following:
% \begin{enumerate}
% \item \textbf{Smooths the Optimization Path}: Momentum reduces the oscillations seen in standard gradient descent, especially along directions with high curvature (e.g., steep walls in valleys).
% \item \textbf{Accelerates Convergence}: By ``remembering'' the gradient history, momentum allows the optimizer to build speed in consistent directions, leading to faster convergence compared to vanilla gradient descent.
% \item \textbf{Overcomes Small Obstacles}: Momentum helps the optimizer overcome small irregularities or plateaus in the loss landscape, avoiding stagnation.
% \end{enumerate}
% 
% Consider optimizing a function with a loss surface resembling a narrow valley.
% Standard gradient descent often zigzags across the valley walls because it reacts directly to the steep slopes.
% Momentum, by contrast, accumulates velocity along the valley's direction,
% leading to smoother and faster movement towards the minimum.
% 
% 
% \subsection{Comparison of Nesterov's Acceleration and Gradient with Momentum}
% 
% Both Nesterov’s Accelerated Gradient (NAG) and Gradient Descent with Momentum (GDM) are optimization techniques 
% that aim to speed up the convergence of gradient-based optimization, particularly for convex and non-convex functions.
% While both methods incorporate momentum to achieve smoother updates, their mechanisms and results differ significantly.
% 
% Gradient Descent with Momentum (GDM) relies on a velocity term that accumulates gradients over iterations. The update rule is:
% \[
% \mathbf{v}_{t+1} = \beta \mathbf{v}_t - \eta \nabla f(\mathbf{x}_t), \quad \mathbf{x}_{t+1} = \mathbf{x}_t + \mathbf{v}_{t+1}.
% \]
% Here, \(\beta\) is the momentum coefficient, and \(\eta\) is the learning rate.
% This method smooths the optimization path, reduces oscillations, and accelerates convergence by incorporating information from past gradients.
% 
% Nesterov’s Accelerated Gradient (NAG), on the other hand, evaluates the gradient at a ``look-ahead'' position,
% anticipating where momentum will take the updates:
% \[
% \mathbf{v}_{t+1} = \beta \mathbf{v}_t - \eta \nabla f(\mathbf{x}_t + \beta \mathbf{v}_t), \quad \mathbf{x}_{t+1} = \mathbf{x}_t + \mathbf{v}_{t+1}.
% \]
% This ``look-ahead'' mechanism allows for more controlled updates by accounting for the influence of momentum beforehand,
% reducing overshooting and improving stability.
% 
% From a convergence perspective, NAG achieves faster rates for convex functions.
% While GDM converges at a rate of \(O(1/k)\) for convex problems, NAG accelerates this to \(O(1/k^2)\).
% This improvement stems from the proactive adjustment of the gradient direction in NAG.
% However, both methods require careful tuning of hyperparameters like the learning rate and momentum coefficient.
% 
% In terms of practical implementation,
% GDM is simpler and widely used due to its straightforward computation of gradients at the current position.
% NAG, while slightly more complex, provides better stability and is less prone to oscillations near optima,
% especially in strongly convex settings.
% 
% Intuitively, GDM can be likened to a ball rolling downhill, gaining momentum as it moves.
% NAG, in contrast, predicts the ball’s future trajectory based on its current velocity and adjusts accordingly,
% leading to smoother and faster convergence.
% 
% In summary, Nesterov’s method is typically preferred when stability and faster convergence are critical,
% particularly for convex problems.
% Gradient Descent with Momentum remains a reliable, simpler alternative,
% especially for non-convex problems or when computational simplicity is prioritized.


\section{Variants of Gradient Descent}

In this section, we compare popular optimization algorithms, including RMSprop, Adagrad, and Adam, against Gradient Descent with Momentum (GDM). These optimizers are widely used for training deep learning models and differ in how they adapt learning rates and incorporate past gradients.

\subsection{Gradient Descent with Momentum (GDM)}
Gradient Descent with Momentum accelerates convergence by adding a fraction of the previous update to the current update. It is particularly effective on loss surfaces with steep gradients or oscillations.

\textbf{Update Rule:}
\[
\mathbf{v}_t = \beta \mathbf{v}_{t-1} + \eta \nabla_{\mathbf{w}} \mathcal{L}(\mathbf{w}_t),
\]
\[
\mathbf{w}_{t+1} = \mathbf{w}_t - \mathbf{v}_t,
\]
where:
\begin{itemize}
    \item $\mathbf{v}_t$ is the velocity (momentum) term.
    \item $\beta \in [0, 1)$ controls the contribution of past gradients.
    \item $\eta$ is the learning rate.
\end{itemize}

Momentum helps the optimizer navigate through ravines and accelerates convergence in directions of consistent gradients while dampening oscillations.

\subsection{RMSprop}
RMSprop (Root Mean Square Propagation) adjusts the learning rate for each parameter based on the magnitude of recent gradients. It normalizes gradients by their running average, preventing excessively large updates and stabilizing training.

\textbf{Update Rule:}
\[
\mathbf{s}_t = \gamma \mathbf{s}_{t-1} + (1 - \gamma) (\nabla_{\mathbf{w}} \mathcal{L}(\mathbf{w}))^2,
\]
\[
\mathbf{w}_{t+1} = \mathbf{w}_t - \frac{\eta}{\sqrt{\mathbf{s}_t + \epsilon}} \nabla_{\mathbf{w}} \mathcal{L}(\mathbf{w}),
\]
where:
\begin{itemize}
    \item The square of Gradient is an element-wise operation
    \item Square root and division $\sqrt{\mathbf{s}_t+\epsilon}$ is also an element-wise operation.
    \item $\mathbf{s}_t$ is the exponentially decaying average of squared gradients.
    \item $\gamma$ is the decay rate, typically set to 0.9.
    \item $\epsilon$ is a small constant added for numerical stability.
\end{itemize}

RMSprop works well for non-stationary objectives and is particularly effective in deep learning tasks.

\subsection{Adam}
Adam (Adaptive Moment Estimation) combines the benefits of RMSprop and Momentum. It maintains a moving average of both the gradients and their squared values, providing an adaptive learning rate that stabilizes updates.

\textbf{Update Rule:}
\[
\mathbf{v}_t = \beta_1 \mathbf{v}_{t-1} + (1 - \beta_1) \nabla_{\mathbf{w}} \mathcal{L}(\mathbf{w}),
\]
\[
\mathbf{s}_t = \beta_2 \mathbf{s}_{t-1} + (1 - \beta_2) (\nabla_{\mathbf{w}} \mathcal{L}(\mathbf{w}))^2,
\]
\[
\hat{\mathbf{v}}_t = \frac{\mathbf{v}_t}{1 - \beta_1^t}, \quad \hat{\mathbf{s}}_t = \frac{\mathbf{s}_t}{1 - \beta_2^t},
\]
\[
\mathbf{w}_{t+1} = \mathbf{w}_t - \frac{\eta}{\sqrt{\hat{\mathbf{s}}_t + \epsilon}} \hat{\mathbf{v}}_t,
\]
where:
\begin{itemize}
    \item $\mathbf{v}_t$ and $\mathbf{s}_t$ are the first and second moment estimates, respectively.
    \item $\beta_1$ and $\beta_2$ control the decay rates for the moments, typically set to 0.9 and 0.999, respectively.
    \item $\epsilon$ ensures numerical stability.
\end{itemize}

Adam is highly effective in practice, offering robust convergence across a variety of tasks and settings.

\subsection{Key Differences}
\begin{itemize}
    \item \textbf{Gradient Adaptation:} RMSprop adapt learning rates for each parameter,
    while GDM applies a single global learning rate with added momentum.
    \item \textbf{Efficiency:} Adam combines the strengths of RMSprop and Momentum, making it a default choice in many deep learning applications.
    \item \textbf{Decay Rates:} RMSprop and Adam rely on exponential moving averages, while Adagrad accumulates gradients over time, leading to different behaviors during training.
\end{itemize}

%%% Sharpness Aware Minimization
%%%%%%%%%%%%%%%%%%%%%%%%%%%%%%%%%%%%%%%%%%%%%%%%%
\section{Sharpness-Aware Minimization (SAM)}

Deep learning models often exhibit overfitting and sensitivity to perturbations in input data or model parameters, primarily due to the sharpness of the loss landscape. Sharp minima, characterized by steep gradients around the optimal parameter values, make models sensitive to small changes, leading to poor generalization on unseen data. This issue is especially pronounced in over-parameterized models, such as deep neural networks, where numerous solutions may perform equally well on the training set but differ significantly in generalization.

To address this, \textbf{Sharpness-Aware Minimization (SAM)} focuses on finding flat minima, defined as regions in the loss landscape where the loss remains relatively constant under small perturbations of the model parameters. Flat minima are empirically correlated with better generalization properties, making SAM an effective method for robust optimization.

Given a loss function $\mathcal{L}(\mathbf{w}, \mathcal{D})$ for model parameters $\mathbf{w}$ over a dataset $\mathcal{D}$, SAM seeks to optimize a surrogate objective that incorporates worst-case perturbations in the parameter space. The objective is defined as:
\[
\min_{\mathbf{w}} \max_{\|\boldsymbol{\epsilon}\|_p \leq \rho} \mathcal{L}(\mathbf{w} + \boldsymbol{\epsilon}, \mathcal{D}),
\]
where $\|\boldsymbol{\epsilon}\|_p \leq \rho$ constrains the perturbation $\boldsymbol{\epsilon}$ to a $p$-norm ball of radius $\rho$. This formulation ensures that the learned model is robust to small perturbations in parameter space, favoring flatter minima.

\textbf{Theorem (Informal):} \emph{Minimizing the sharpness-aware objective leads to a flatter solution in the loss landscape, which empirically correlates with better generalization performance.}

\noindent This theorem encapsulates the core intuition of SAM. By optimizing for a worst-case adversarial perturbation, the solution lies in a flat region of the loss landscape. Flatness, a geometric property of the loss landscape, is strongly associated with model robustness and improved generalization.

To solve the min-max optimization problem efficiently, SAM introduces a two-step iterative process:

\begin{enumerate}
    \item \textbf{Gradient Ascent (Inner Maximization):} Compute the adversarial perturbation $\boldsymbol{\epsilon}^*$ that maximizes the loss in a neighborhood of the current parameters $\mathbf{w}$. Using a first-order approximation and the $\ell_2$-norm, this is given by:
    \[
    \boldsymbol{\epsilon}^* = \rho \frac{\nabla_{\mathbf{w}} \mathcal{L}(\mathbf{w}, \mathcal{D})}{\|\nabla_{\mathbf{w}} \mathcal{L}(\mathbf{w}, \mathcal{D})\|_2}.
    \]

    \item \textbf{Gradient Descent (Outer Minimization):} Update the model parameters $\mathbf{w}$ to minimize the adversarially perturbed loss:
    \[
    \mathbf{w} \leftarrow \mathbf{w} - \eta \nabla_{\mathbf{w}} \mathcal{L}(\mathbf{w} + \boldsymbol{\epsilon}^*, \mathcal{D}),
    \]
    where $\eta$ is the learning rate.
\end{enumerate}

The additional adversarial perturbation introduced by SAM ensures that the optimization process favors flatter minima in the loss landscape, which are empirically linked to improved generalization. The computational overhead of SAM is manageable, requiring only one additional gradient computation per iteration. This makes SAM a practical and effective approach for robust deep learning optimization.



% \begin{theorem}
%     Suppose $f$ is a convex function where $\nabla f$ is $L$-Lipschitz.
%     Let learning rate be $\alpha_k\equiv \alpha = 1/L$, the Gradient Descent algorithm converges in rate $O(1/k)$.
% \end{theorem}
% \begin{proof}
%     Let $x^\star$ be a global optimum.
%     From the previous theorem, GD can achieve $x^\star$.
%     With $\alpha_k \equiv \alpha$ (for simplicity),
%     \begin{align}
%         f(x^{(k+1)}) - f(x^{(0)})
%         &= \sum_{i=0}^k f(x^{(i+1)}) - f(x^{(i)})\\
%         \leq & - \frac{1}{2L} \sum_{i=0}^k \|\nabla f(x^{(i)})\|^2
%     \end{align} 
%     Then, 
% \end{proof}
% 
% If strong convex,
% 
% \noa{meaning of strong convex}



\bibliographystyle{abbrv}
\bibliography{scribe-ref}

\end{document}